% !TeX root = ../../Book3.tex
% 纲要标题：## 第26章*　抛物方法与调和映射热流
% 纲要位置：几何学计划纲要.md，第 3198 行。
% 正文状态：初稿；尚未正式审核。

% 纲要内容开始
% 
% > 本章目标与证明范围：作为几何流路线的工具入口，可不经过 Yamabe 和 Einstein 章。先证明紧光滑情形所需线性及拟线性短时理论，再在非正目标曲率下完成 Eells–Sampson 型代表模块。 非线性共同输入定位于 III.13.5—13.6，不以完整阅读 Yamabe 章为先修。
% 
% 纲要内容结束

\OptionalChapter{抛物方法与调和映射热流}{g3:ch:26}
\BookChapterNumber{g3:ch:26}

\begin{chapterabstract}
	用能量下降的演化寻找调和映射，首先需要一个可证明、可延拓的抛物初值理论。本章从热核与 Duhamel 公式出发，在抛物能量空间中证明一般强抛物线性系统的初值同构、正则性和标量最大值原理；再以定量求逆构造拟线性短时解，证明唯一性、光滑依赖以及基于二阶导数控制的延拓判据。对调和映射热流，逐项计算总能量、梯度和时间速度的演化。在闭源与紧非正曲率目标之间，固定长度时间窗上的热核估计给出一致梯度界，协变导数演化进一步给出正时间后的一致高阶界，从而得到长期存在及同伦类中的调和代表。最后以测地同伦的能量凸性讨论极小性、唯一性和附加条件下的全时收敛。
\end{chapterabstract}

\begin{learninggoals}
	\begin{itemize}
		\item 用热半群解释非齐次方程，并确定抛物能量空间中的初值和时间迹。
		\item 从强抛物主符号推导能量估计，通过有限维逼近证明线性存在与唯一性。
		\item 验证拟线性方程的 Banach 空间映射性质，构造短时解并证明光滑依赖与延拓判据。
		\item 在非负 Laplace 约定下推导调和映射热流的能量、梯度及时间速度演化。
		\item 把固定时间的热传播、能量下降和高阶导数估计结合，证明非正曲率目标下的全时存在。
		\item 由光滑紧性构造保持同伦类的调和极限，并辨别子列收敛、全时收敛和调和代表唯一性。
	\end{itemize}
\end{learninggoals}

在圆周之间考虑映射 \(u_0([x])=[dx+a\sin x]\)，其中 \(d\) 是整数。线性部分记录映射绕目标圆周的次数，周期振荡则增加能量。把它放入热方程，得到 \(u([x],t)=[dx+a e^{-t}\sin x]\)：振荡逐渐消失，次数保持不变，最后留下一个调和映射。这个计算提示了一种几何构造：沿能量下降方向移动映射，在保留同伦信息的同时消去多余的空间振荡。

\ProseParagraphBreak
一般目标上，下降方向是张力场 \(\tau(u)\)，相应方程为 \(\partial_tu=\tau(u)\)。局部坐标中的最高阶部分仍是热算子，但非线性项含目标联络和一阶导数的二次乘积。要把圆周算例的方法用于一般流形，必须先证明这个方程有解，再防止梯度与高阶导数在有限时间失控。能量下降只是一个积分控制，不能单独排除这些问题。

\ProseParagraphBreak
非正目标曲率使梯度与时间速度的 Bochner 演化具有有利符号。源的 Ricci 曲率可以有负的部分，因此直接应用最大值原理首先得到随时间增长的梯度上界；将总能量界与长度固定的热传播时间结合，才能得到不依赖终止时间的一致上界。随后以协变导数演化逐阶控制空间振荡，便可在无穷时间中抽取光滑极限。Eells--Sampson 的调和代表定理正体现了这条思路。\footnote{James Eells, Jr.（1926--2007），美国数学家，研究整体分析与微分几何。Marco Joseph Harold Sampson, Jr.（1925--2003），生于美国的数学家，研究分析、几何与调和映射；其生年据\href{https://aspace.library.jhu.edu/repositories/3/resources/801}{Johns Hopkins 大学档案}。}原文的热核、逐次逼近及长期存在可参看\cite[第9--10节，第145--155页]{EellsSampson1964HarmonicMaps}，本章使用抛物能量空间组织短时证明。

\ProseParagraphBreak
所需分析基础是第11—13章的 Sobolev、椭圆及 Banach 空间理论，以及第16章的热核；映射几何来自第18.1节的张力场与变分公式。最后的唯一性证明还用第8章的完备覆盖与第9章的非正曲率测地线性质。前两节同时为后续 Ricci 流与平均曲率流准备短时工具，阅读本章无需先完成 Yamabe 或 Einstein 章。

% !TeX root = ../../Book3.tex
% 纲要标题：### 26.1 线性抛物方程
% 纲要位置：几何学计划纲要.md，第 3202 行。
% 正文状态：初稿；尚未正式审核。

% 纲要内容开始
% 
% * 热核和 Duhamel 公式
% * 抛物 Hölder 或 Sobolev 空间
% * 线性估计与初值相容
% * 抛物最大值原理
% 
% 纲要内容结束

\section{线性抛物方程}
\label{g3:sec:2601}

热方程同时描述两件事：初值沿时间向前传播，空间振荡逐渐衰减。第十六章已经用谱展开构造了固定 Laplace 算子的热核。几何流还会产生随时间和位置变化的系数，因此需要一种不依赖显式热核的求解方法。本节用能量估计和有限维逼近建立这个方法，并把初值所处的空间一并确定。

\ProseParagraphBreak
本节的空间流形始终闭，即紧且无边界；积分使用固定黎曼度量的体积密度。先讨论标量热方程，再允许固定光滑向量丛上的矩阵主部。标量函数的大小次序将在最后一小节用于最大值原理。

\subsection{热核和 Duhamel 公式}
\label{g3:subsec:2601:01}
设 \((M,g)\) 为闭光滑黎曼流形，\(\dim M=n\ge1\)，仍取 \(\Delta=-\operatorname{div}_g\operatorname{grad}_g\)。\Cref{g3:prop:heat-semigroup-properties,g3:thm:heat-spectral-solution,g3:thm:heat-kernel-spectral}已经证明：\(T_t=e^{-t\Delta}\) 是强连续半群，在正时间平滑初值，并由热核 \(K(t,x,y)\) 表示。非齐次项可以看成不断加入的新初值；在时间 \(a\) 加入的部分，到时间 \(t\) 已传播了 \(t-a\)。

\ProseParagraphBreak
将这些传播后的贡献相加，就得到\term[Duhamel gongshi]{Duhamel 公式}[Duhamel's formula]。\footnote{Jean-Marie Constant Duhamel（1797--1872），法国数学家、物理学家，研究偏微分方程及热传导、力学与声学。}

\begin{proposition}[Duhamel 公式]\label{g3:prop:parabolic-duhamel}
	设 \((M,g)\) 为正维闭光滑黎曼流形，\(T>0\)，\(u_0\in C^\infty(M)\)、\(f\in C^\infty([0,T]\times M)\)。则
	\begin{equation}\label{g3:eq:parabolic-duhamel}
		u(t)=T_tu_0+\int_0^tT_{t-a}f(a)\,\mathrm da
	\end{equation}
	是 \(\partial_tu+\Delta u=f\)、\(u(0)=u_0\) 的唯一光滑解，光滑性在时间端点按单侧导数理解。对 \(t>0\)，同一公式可写为
	\[
		u(t,x)=\int_MK(t,x,y)u_0(y)\,\mu_g(y)
		+\int_0^t\int_MK(t-a,x,y)f(a,y)\,\mu_g(y)\,\mathrm da.
	\]
\end{proposition}
\begin{proof}
	取第十六章的正交特征基 \(\Delta\phi_j=\lambda_j\phi_j\)，记 \(f_j(t)=(f(t),\phi_j)\)、\(u_{0j}=(u_0,\phi_j)\)。对方程取第 \(j\) 个系数，得到
	\[
		u_j'+\lambda_ju_j=f_j,
		\qquad
		u_j(t)=e^{-t\lambda_j}u_{0j}
		+\int_0^te^{-(t-a)\lambda_j}f_j(a)\,\mathrm da.
	\]
	这是乘以积分因子后直接积分所得。由于 \(f\) 光滑，其每个固定空间 Sobolev 范数在 \([0,T]\) 上有界。热半群在相应谱范数中压缩，故上式的级数和时间积分在每个这样的范数中收敛。先用光滑有限谱和求导，得 \(u_t=f-\Delta u\)；右端在任意固定 Sobolev 范数中连续。再对时间求导，并逐次提高所用空间阶数，得到所有混合导数。\Cref{g3:cor:higher-sobolev-embedding}于是给出所述联合光滑性和初值。

	任意另一光滑解的各个特征系数满足同一个标量初值问题，所以逐项相同，完备性给出唯一性。最后代入已经证明的热核表示。时间积分可先在 \([0,t-\varepsilon]\) 上写成核积分，再令 \(\varepsilon\downarrow0\)；余下时间段的任意固定 Sobolev 范数至多为 \(\varepsilon\sup_a\|f(a)\|_{H^k}\)，故极限确实等于式\eqref{g3:eq:parabolic-duhamel}。
\end{proof}

\begin{example}\label{g3:exm:heat-forcing-zero-mode}
	在连通闭流形上，若外力不随时间变化，写 \(f=\sum_jb_j\phi_j\)，且初值为零，则
	\[
		u(t)=t\sum_{\lambda_j=0}b_j\phi_j
		+\sum_{\lambda_j>0}\frac{1-e^{-t\lambda_j}}{\lambda_j}b_j\phi_j.
	\]
	非零平均值持续积累成线性增长；正频率则趋向相应的 Poisson 解分量。例如圆周上的外力 \(1+\cos x\) 给出 \(u(t,x)=t+(1-e^{-t})\cos x\)。因此正时间平滑与长时间有界是两个不同的问题。
\end{example}

对仅有平方可积正则性的数据，式\eqref{g3:eq:parabolic-duhamel}仍有积分意义，但求导及时间端点的强收敛需要估计。下面的空间同时记录这两项要求，也适用于没有固定热半群的时变系统。

\subsection{抛物 Hölder 或 Sobolev 空间}
\label{g3:subsec:2601:02}
固定闭光滑黎曼流形 \((M^n,g)\)、\(n\ge1\)，以及正秩光滑度量向量丛 \(E\to M\) 和度量联络 \(\nabla\)。丛可以为实或复；复情形的内积对第一变量线性。空间 Sobolev 范数取自这些固定数据。负一阶空间 \(H^{-1}(E)\) 是 \(H^1(E)\) 关于 \(L^2\) 配对的对偶。

\ProseParagraphBreak
为了在逼近中同时控制不同导数阶数，令 \(P=I+\nabla^*\nabla\)。\Cref{g3:thm:positive-elliptic-realization,g3:thm:elliptic-discrete-spectrum}给出光滑正交特征基 \((e_j)_{j\ge1}\)，满足 \(Pe_j=\mu_je_j\)、\(\mu_j\ge1\)。实丛可先复化；因 \(P\) 保持实截面，每个特征空间都可取实正交基。对整数 \(r\ge0\)，可以使用等价范数
\begin{equation}\label{g3:eq:parabolic-spectral-norm}
	\|v\|_r^2=\sum_j\mu_j^r|v_j|^2,
	\qquad v_j=(v,e_j)_{L^2}.
\end{equation}
偶数阶的等价性由\cref{g3:thm:elliptic-spectral-sobolev}给出。若 \(r=2q+1\)，则对光滑有限和有
\[
	\sum_j\mu_j^{2q+1}|v_j|^2
	=\|P^qv\|_{L^2}^2+\|\nabla P^qv\|_{L^2}^2.
\]
右侧至多为 \(C\|v\|_{H^{2q+1}}^2\)；反向对 \(P^q\) 使用\cref{g3:thm:global-elliptic-hk}的高一阶估计，并用 \(\|v\|_{L^2}\le\|P^qv\|_{L^2}\)，得到同一阶范数的控制。\(q=0\) 时就是 \(H^1\) 范数的定义。由光滑稠密性和有限谱和的完备化，结论对全部 \(H^r\) 截面成立。对 \(H^{-1}\) 使用权 \(\mu_j^{-1}\)，与上面的对偶定义一致。

\ProseParagraphBreak
记 \(\Pi_N\) 为前 \(N\) 个特征截面的正交投影。它在每个这些范数中都是压缩，且强收敛到恒等。对 \(a\in H^{r-1}\)、\(b\in H^{r+1}\)，定义配对
\[
	[a,b]_r=\sum_j\mu_j^ra_j\overline{b_j},
	\qquad |[a,b]_r|\le\|a\|_{r-1}\|b\|_{r+1}.
\]
当两者光滑时，这等于 \((a,P^rb)_{L^2}\)。

\begin{lemma}\label{g3:lem:parabolic-energy-trace}
	设 \(E\) 为闭光滑黎曼流形上的光滑度量丛，\(T>0\)、整数 \(r\ge0\)。若
	\[
		u\in L^2(0,T;H^{r+1}(E)),\qquad
		\partial_tu\in L^2(0,T;H^{r-1}(E))
	\]
	且时间导数按分布理解，则 \(u\) 有唯一的 \(C([0,T];H^r(E))\) 代表。对 \(0\le a\le b\le T\)，有
	\begin{equation}\label{g3:eq:parabolic-energy-chain}
		\|u(b)\|_r^2-\|u(a)\|_r^2
		=2\operatorname{Re}\int_a^b[\partial_tu(t),u(t)]_r\,\mathrm dt.
	\end{equation}
	时间迹 \(u\mapsto u(0)\) 对上述两个 \(L^2\) 范数之和连续；这一迹估计的常数允许依赖 \(T\)。
\end{lemma}
\begin{proof}
	每个系数 \(u_j(t)\) 及其弱导数属于 \(L^2(0,T)\)，所以有绝对连续代表。一维结论可由积分 \(u_j'\) 并使用“分布导数为零则为常数”直接得到。因此有限和 \(\Pi_Nu\) 连续，且满足式\eqref{g3:eq:parabolic-energy-chain}的有限维版本。

	对 \(N>m\)，令 \(w=(\Pi_N-\Pi_m)u\)。存在 \(a\in[0,T]\)，使 \(\|w(a)\|_r^2\le T^{-1}\|w\|_{L^2H^r}^2\)。从这个时刻向两侧积分有限维能量恒等式，得到
	\[
		\sup_{0\le t\le T}\|w(t)\|_r^2
		\le T^{-1}\|w\|_{L^2H^r}^2
		+2\|w_t\|_{L^2H^{r-1}}\|w\|_{L^2H^{r+1}}.
	\]
	两种谱尾项都趋于零，故 \(\Pi_Nu\) 在 \(C([0,T];H^r)\) 中为 Cauchy 列。极限与原来的 \(L^2\) 截面相同，给出所需代表。对有限和的能量等式取极限，右侧用两种 \(L^2\) 强收敛及配对不等式，得到式\eqref{g3:eq:parabolic-energy-chain}。上述估计也对完整的 \(u\) 成立，因而给出迹的连续性。两个连续代表若几乎处处相同，便处处相同。
\end{proof}

\begin{definition}\label{g3:def:parabolic-energy-space}
	在上述固定几何数据下，\(T>0\)、整数 \(r\ge0\) 时，记
	\begin{equation}\label{g3:eq:parabolic-energy-space}
		\begin{split}
		X_T^r(E)=\{u:\;&u\in C([0,T];H^r(E))\cap L^2(0,T;H^{r+1}(E)),\\
		&\partial_tu\in L^2(0,T;H^{r-1}(E))\}.
		\end{split}
	\end{equation}
	赋予它三个所列范数之和，称为本章的\term[paowu nengliang kongjian]{抛物能量空间}[parabolic energy space]。其零初值闭子空间记为 \(X_{T,0}^r(E)\)。
\end{definition}

完备性可逐项核对：Cauchy 列在连续函数空间和两种 \(L^2\) 空间中各有极限，测试配对说明极限的弱时间导数仍是导数项的极限。上一引理则说明连续时间代表不是额外的光滑假设。这里空间增益以 \(L^2\) 时间积分记录：初值仅在 \(H^r\) 中，正时间的 \(H^{r+1}\) 范数却平方可积。

\begin{proposition}\label{g3:prop:parabolic-zero-trace-small-time}
	设 \(M^n\) 闭且光滑，\(E\) 为固定光滑度量丛，整数 \(s>n/2+3\)。若 \(u\in X_{T,0}^s(E)\)，则
	\[
		\sup_{0\le t\le T}\|u(t)\|_{H^{s-1}}
		\le T^{1/2}\|\partial_tu\|_{L^2H^{s-1}},
		\qquad
		\sup_{0\le t\le T}\|u(t)\|_{C^2}
		\le C T^{1/2}\|u\|_{X_T^s},
	\]
	其中 \(C\) 只依赖固定几何数据和 \(s\)。
\end{proposition}
\begin{proof}
	在 \(H^{s-1}\) 中有 \(u(t)=\int_0^t u_t(a)\,\mathrm da\)，由 Cauchy--Schwarz 得第一式。因 \(s-1>n/2+2\)，\cref{g3:cor:higher-sobolev-embedding}将它转为第二式。
\end{proof}

这给短时非线性论证一个有用的量：零初值修正虽然未必在最高阶连续 Sobolev 范数中自动变小，其二阶连续范数却带有 \(T^{1/2}\)。

\subsection{线性估计与初值相容}
\label{g3:subsec:2601:03}
现在让系数依赖 \((t,x)\)。固定 \(T_0>0\)，在同一闭流形、丛及联络上考虑
\begin{equation}\label{g3:eq:linear-parabolic-operator}
	L(t)u=\nabla^*(A(t)\nabla u)+B(t)\nabla u+C(t)u.
\end{equation}
这里 \(A\) 是 \(\mathrm T^*M\otimes E\) 的丛自同态，\(B:\mathrm T^*M\otimes E\to E\)、\(C:E\to E\)，均在 \([0,T_0]\times M\) 上光滑。在正交标架中，主符号写为 \(A^{ij}(t,x)\xi_i\xi_j\)。我们要求存在 \(\theta>0\)，使
\begin{equation}\label{g3:eq:linear-strong-parabolicity}
	\operatorname{Re}\langle A^{ij}(t,x)\xi_i\xi_jv,v\rangle
	\ge\theta|\xi|_g^2|v|^2
	\quad(t\in[0,T_0],\ x\in M,\ \xi\in\mathrm T_x^*M,\ v\in E_x).
\end{equation}
这是本节所用的一致强抛物条件。它允许分量间耦合，也不要求 \(L(t)\) 自伴。任意具有此类主符号的光滑二阶算子，都可把主部写成式\eqref{g3:eq:linear-parabolic-operator}的散度形式，把新增的一阶项并入 \(B\)。

\ProseParagraphBreak
下面的估计允许用低阶范数补偿主部能量的低阶误差，称为\term[Garding budengshi]{Gårding 不等式}[Gårding inequality]。\footnote{Lars Gårding（1919--2014），瑞典数学家，研究偏微分方程与偏微分算子。}

\begin{lemma}\label{g3:lem:parabolic-garding}
	设 \((M^n,g)\) 为正维闭光滑黎曼流形，\(E\) 为正秩光滑度量丛，\(L(t)\) 满足式\eqref{g3:eq:linear-parabolic-operator}及\eqref{g3:eq:linear-strong-parabolicity}，系数在 \([0,T_0]\times M\) 上光滑。对每个整数 \(r\ge0\)，存在 \(c_r>0\)、\(C_r<\infty\)，使全部光滑截面 \(v\) 及全部 \(t\in[0,T_0]\) 满足
	\begin{equation}\label{g3:eq:parabolic-garding}
		\operatorname{Re}(L(t)v,P^rv)_{L^2}
		\ge c_r\|v\|_{r+1}^2-C_r\|v\|_r^2.
	\end{equation}
	此外 \(\|L(t)v\|_{r-1}\le C_r\|v\|_{r+1}\)。常数只需固定几何数据、\(\theta\) 以及系数的有限阶空间导数的一致界，不需时间导数界。
\end{lemma}
\begin{proof}
	先说明为何秩一方向上的条件足以控制积分能量。把系数和背景度量冻结在一个坐标点，且令 \(v\) 紧支于该坐标片。Plancherel 恒等式把最高阶二次型化为
	\[
		\int_{\mathbb R^n}|\xi|^{2r}
		\operatorname{Re}\langle A^{ij}\xi_i\xi_j\widehat v(\xi),\widehat v(\xi)\rangle\,\mathrm d\xi
		\ge\theta\int_{\mathbb R^n}|\xi|^{2r+2}|\widehat v(\xi)|^2\,\mathrm d\xi.
	\]
	在非正交坐标中以冻结的正定二次型代替 \(|\xi|^2\)，范数等价给出相同结论。这里 \(P^r\) 是阶 \(2r\) 的微分算子，主符号为 \(|\xi|^{2r}I\)；分部积分 \(r+1\) 次，将上述最高阶部分写成两边各有 \(r+1\) 阶导数的双线性型。

	缩小坐标片，使主系数与冻结值之差不超过一个给定的 \(\varepsilon\)。冻结误差于是至多为 \(C\varepsilon\|v\|_{H^{r+1}}^2\)。导数落到系数、联络或体积密度时，总共作用在两个 \(v\) 上的阶数至多为 \(2r+1\)；继续分部积分可使每一边至多有 \(r+1\) 阶，并且至少一边至多有 \(r\) 阶。因此这些项由
	\[
		C\|v\|_{H^{r+1}}\|v\|_{H^r}
		\le\varepsilon'\|v\|_{H^{r+1}}^2+C_{\varepsilon'}\|v\|_{H^r}^2
	\]
	控制。\(B,C\) 项具有同样的低阶性质。先选小片和 \(\varepsilon'\)，便能吸收两类最高阶误差。

	取满足 \(\sum_\alpha\chi_\alpha^2=1\) 的有限光滑分割，将局部估计用于 \(\chi_\alpha v\)。比较局部式之和与全局的 \((Lv,P^rv)\) 时，所有额外项都至少有一个导数落在 \(\chi_\alpha\) 上，仍按刚才的低阶估计处理。加上 \(\|v\|_{H^r}^2\) 补齐最高阶齐次范数，再使用式\eqref{g3:eq:parabolic-spectral-norm}，即得式\eqref{g3:eq:parabolic-garding}。\([0,T_0]\times M\) 紧且系数光滑，有限空间覆盖、冻结半径和所用系数导数界可对时间统一选取。

	最后，当 \(r\ge1\) 时，Leibniz 公式直接给出 \(L:H^{r+1}\to H^{r-1}\) 的界。\(r=0\) 时用散度形式与 \(H^1\) 测试截面配对：主项为 \((A\nabla v,\nabla w)\)，低阶项同样受 \(C\|v\|_{H^1}\|w\|_{H^1}\) 控制。这正是到 \(H^{-1}\) 的所需界。
\end{proof}

\begin{theorem}\label{g3:thm:linear-parabolic-isomorphism}
	设 \((M^n,g)\) 为正维闭光滑黎曼流形，\(E\) 为正秩光滑度量丛，\(L(t)\) 满足上一引理的全部条件。对整数 \(r\ge0\)、\(0<T\le T_0\)，映射
	\[
		u\longmapsto(\partial_tu+L(t)u,\ u(0))
	\]
	是从 \(X_T^r(E)\) 到 \(L^2(0,T;H^{r-1}(E))\times H^r(E)\) 的有界线性同构。具体地，对任意 \(u_0\in H^r(E)\) 与 \(f\in L^2(0,T;H^{r-1}(E))\)，唯一解满足
	\begin{equation}\label{g3:eq:linear-parabolic-energy}
		\begin{split}
		&\sup_{0\le t\le T}\|u(t)\|_r^2
		+\int_0^T\bigl(\|u(t)\|_{r+1}^2+\|\partial_tu(t)\|_{r-1}^2\bigr)\,\mathrm dt\\
		&\qquad\le C\left(\|u_0\|_r^2+\int_0^T\|f(t)\|_{r-1}^2\,\mathrm dt\right).
		\end{split}
	\end{equation}
	其中 \(C\) 可对 \(0<T\le T_0\) 一致选取，其依赖为上一引理中的数据及 \(T_0\)。特别地，零初值时 \(\partial_t+L:X_{T,0}^r\to L^2H^{r-1}\) 的逆范数在这些时间区间上一致有界。
\end{theorem}
\begin{proof}
	证明分为能量控制、有限维求解及取极限。关键是选择与用于能量的 \(P^r\) 交换的谱投影，这样求解截断方程不会产生未控制的投影交换子。

	先对光滑解与 \(P^ru\) 配对。由上一引理及 Young 不等式，
	\[
		\begin{aligned}
		\frac12\frac{\mathrm d}{\mathrm dt}\|u\|_r^2
		+c_r\|u\|_{r+1}^2
		&\le C_r\|u\|_r^2+|[f,u]_r|\\
		&\le C_r\|u\|_r^2+\frac{c_r}{2}\|u\|_{r+1}^2+C\|f\|_{r-1}^2.
		\end{aligned}
	\]
	乘以相应的 \(e^{-Ct}\) 并积分，得到 \(\sup_t\|u(t)\|_r^2\) 的控制；再对未乘指数的式子积分，得到 \(L^2H^{r+1}\) 控制。方程和 \(L:H^{r+1}\to H^{r-1}\) 的界又控制 \(u_t\)。这给出式\eqref{g3:eq:linear-parabolic-energy}，且所有指数因子均可用 \(e^{CT_0}\) 控制。

	令 \(u_N(t)\) 取值于 \(\operatorname{span}\{e_1,\ldots,e_N\}\)，求解
	\[
		\partial_tu_N+\Pi_NL(t)u_N=\Pi_Nf,
		\qquad u_N(0)=\Pi_Nu_0.
	\]
	这是系数连续、右端平方可积的有限维线性常微分方程。其基本解矩阵可由积分方程的逐次逼近构造，再用变常数公式得到绝对连续解；有限时间上的系数有界保证解延拓至 \(T\)。由于 \(\Pi_NP^ru_N=P^ru_N\)，刚才的能量计算逐字适用于 \(u_N\)。同时 \(\Pi_N\) 在 \(H^{r-1}\) 中压缩，所以时间导数的界也与 \(N\) 无关。

	因此可取子列，使 \(u_N\) 在 \(L^2H^{r+1}\) 中弱收敛到 \(u\)，而 \(\partial_tu_N\) 在 \(L^2H^{r-1}\) 中弱收敛。对光滑时间测试函数和有限谱截面配对，积分分部说明后一极限就是 \(\partial_tu\)。同样的配对允许在有限维方程中取极限，得到 \(u_t+L(t)u=f\)：\(L(t)\) 是这两种空间之间一致有界的线性算子，有限谱测试又在空间测试类中稠密。

	\Cref{g3:lem:parabolic-energy-trace}使 \(u\) 有连续 \(H^r\) 代表，并使时间迹成为上述 Hilbert 交空间上的连续线性映射。故 \(u_N(0)\rightharpoonup u(0)\) 于 \(H^r\)，而 \(\Pi_Nu_0\to u_0\)，从而 \(u(0)=u_0\)。对极限解使用式\eqref{g3:eq:parabolic-energy-chain}，其余配对由 \(H^{r+1}\) 中的光滑稠密性成立，于是同一能量估计适用于 \(u\)。

	若两个解的数据相同，其差的数据为零，式\eqref{g3:eq:linear-parabolic-energy}迫使差恒为零。一般两个数据之差满足同一估计，因此求解映射线性且有界。正向映射的有界性由空间定义、时间迹及 \(L\) 的有界性给出，证明了有界同构。
\end{proof}

\begin{theorem}\label{g3:thm:linear-parabolic-smoothing}
	设 \(M,E,L\) 满足\cref{g3:thm:linear-parabolic-isomorphism}的条件，\(0<T\le T_0\)，\(r\ge0\) 为整数。若 \(u_0\in C^\infty(E)\)、\(f\in C^\infty([0,T]\times M;E)\)，则相应解在 \([0,T]\times M\) 上光滑。若只假定 \(u_0\in H^r\)、\(f\in L^2(0,T;H^{r-1})\)，而 \(f\) 在 \((0,T]\times M\) 上光滑，则解在 \((0,T]\times M\) 上光滑。
	更具体地，若整数 \(k\ge1\)、\(0<\tau<T\le T_0\)，且 \(f\in L^2(\tau/2,T;H^{r+k-1})\)，则
	\begin{equation}\label{g3:eq:linear-parabolic-positive-time}
		\|u\|_{X^{r+k}_{[\tau,T]}}
		\le C_{r,k,\tau}\left(
		\|u_0\|_r+\|f\|_{L^2(0,T;H^{r-1})}
		+\|f\|_{L^2(\tau/2,T;H^{r+k-1})}\right),
	\end{equation}
	其中区间版本的 \(X\) 范数按式\eqref{g3:eq:parabolic-energy-space}定义，常数还依赖固定几何数据、\(T_0\) 和所用系数界。
\end{theorem}
\begin{proof}
	光滑数据属于每个整数阶的数据空间。分别对所有 \(r\) 应用同构定理，并用零阶唯一性识别这些解，就得到 \(u\in C([0,T];H^m)\) 对所有 \(m\) 成立。方程给出 \(u_t=f-L(t)u\in C([0,T];H^m)\) 对所有 \(m\) 成立。继续求时间导数，右端只含系数的时间导数和已经得到的较低时间导数，归纳可得任意阶。空间 Sobolev 嵌入将它们转为联合光滑性。

	再考虑较粗初值。能量估计给出 \(u\in L^2(0,T;H^{r+1})\)。在任何正长度的子区间中，可选时刻 \(a\)，使 \(u(a)\in H^{r+1}\)，且这一范数不超过该区间上 \(L^2H^{r+1}\) 范数除以区间长度的平方根。从 \(a\) 出发应用高一阶的同构定理，所得解由唯一性与原解一致，故获得下一阶能量控制。

	为得到式\eqref{g3:eq:linear-parabolic-positive-time}，将 \((\tau/2,\tau)\) 分成 \(k\) 个正长度小区间，每次在下一个小区间选择上述时刻。第 \(j\) 次用 \(r+j\) 阶的同构定理；所需外力范数受 \(L^2(\tau/2,T;H^{r+k-1})\) 控制。经过 \(k\) 次，便在 \([\tau,T]\) 上得到所列估计。所有长度和各次能量常数都由 \(k,\tau\) 及所列数据控制。若外力正时间光滑，可对任意 \(k,\tau\) 使用这一结论，再由方程取得时间导数，证明正时间光滑性。
\end{proof}

由于没有空间边界，初值不必另外满足边界相容条件。若希望解光滑到 \(t=0\)，它的时间导数却已被方程确定：记 \(u_j=\partial_t^ju|_{t=0}\)，则
\begin{equation}\label{g3:eq:parabolic-initial-jets}
	 u_{j+1}=\partial_t^jf(0)
	-\sum_{\ell=0}^j\binom j\ell(\partial_t^\ell L)(0)u_{j-\ell}.
\end{equation}
这是一组由数据算出的初始喷射，而不是可以自由再指定的初值。带空间边界的问题还须将这些喷射代入相应边界条件；本章的闭流形结论不包含那一额外问题。

\ProseParagraphBreak
当 \(L=\Delta\) 时，式\eqref{g3:eq:parabolic-duhamel}与这里的能量解一致。对 \(u_0\in H^r\)、\(f\in L^2H^{r-1}\)，先取光滑数据逼近，能量估计给出 \(X_T^r\) 收敛；半群在 \(H^{r-1}\) 中有界，时间积分又由 \(L^2\) 控制为 \(L^1\)，所以 Duhamel 公式可在 \(H^{r-1}\) 中取极限。其值实际上具有已经证明的连续 \(H^r\) 代表。

\subsection{抛物最大值原理}
\label{g3:subsec:2601:04}
能量控制一个空间积分，最大值原理则直接控制每个时刻的最大值。后者利用实函数在极大点的一阶导数为零、Hessian 非正这一事实，因此先限于标量方程。以下以 \(\nabla^2\) 表示固定背景度量的 Hessian。

\begin{theorem}\label{g3:thm:parabolic-maximum}
	设 \((M^n,g)\) 为正维闭光滑黎曼流形，\(T>0\)，实函数 \(u\) 连续于 \([0,T]\times M\)，并在正时间具有连续的一阶时间、二阶空间导数。设实系数 \(a^{ij},b^i,c\) 及实函数 \(f\) 在 \([0,T]\times M\) 上连续，\(a^{ij}\) 对称且一致正定，\(c\ge-\gamma\)、\(\gamma\ge0\)。若
	\[
		\partial_tu-a^{ij}\nabla_i\nabla_ju+b^i\nabla_i u+cu\le f,
		\qquad f(t,x)\le F(t),
	\]
	其中 \(F\) 为非负连续函数，则
	\begin{equation}\label{g3:eq:parabolic-maximum-bound}
		u(t,x)\le e^{\gamma t}\max\{0,\max_Mu(0)\}
		+\int_0^te^{\gamma(t-a)}F(a)\,\mathrm da.
	\end{equation}
	特别地，若 \(u\) 满足对应等式，则将 \(F(t)\) 取为 \(\|f(t)\|_\infty\)，有相同形式的 \(\|u(t)\|_\infty\) 估计，初项改为 \(\|u(0)\|_\infty\)。
\end{theorem}
\begin{proof}
	令式\eqref{g3:eq:parabolic-maximum-bound}的右端为 \(m(t)\)。则 \(m\ge0\)、\(m'=\gamma m+F\)，所以 \(m'+cm\ge F\)。记 \(w=u-m\)，便有 \(w(0)\le0\)，且 \(w\) 满足同一算子的不等式、右端为零。

	令 \(v=e^{-\gamma t}w\)，则
	\[
		\partial_tv-a^{ij}\nabla_i\nabla_jv+b^i\nabla_iv+(c+\gamma)v\le0,
	\]
	其中零阶系数非负。对 \(\varepsilon>0\)，取 \(v_\varepsilon=v-\varepsilon e^t\)，其左端严格小于零，初值严格负。若它在某个 \([0,t_1]\times M\) 上有正最大值，最大值点的时间必大于零；在那里空间梯度为零、Hessian 非正，时间导数非负。即使最大值在上端时刻 \(t_1\) 取得，左侧时间差商仍给出该符号。因此算子在此点的值非负，与严格不等式矛盾。故 \(v_\varepsilon\le0\)，令 \(\varepsilon\downarrow0\) 得 \(w\le0\)。对等式中的 \(u\) 和 \(-u\) 分别应用结论，即得双侧估计。
\end{proof}

\begin{corollary}\label{g3:cor:parabolic-gronwall}
	设 \((M,g)\) 为正维闭光滑黎曼流形，\(v\ge0\) 连续到初时刻并在正时间属于 \(C^{1,2}\)。若 \(a\ge0\) 为常数且
	\[
		(\partial_t+\Delta_g)v\le av,
	\]
	则 \(\max_Mv(t)\le e^{at}\max_Mv(0)\)。若两个实标量函数的差满足上一最大值定理中右端为零的不等式，且初值之差非正，则以后仍保持这个次序。
\end{corollary}
\begin{proof}
	第一项取 \(c=-a\)、\(F=0\)，第二项直接将最大值定理用于两函数之差。
\end{proof}

\begin{example}\label{g3:exm:parabolic-system-order-failure}
	在任意闭流形上取平凡二维实丛，考虑
	\[
		\partial_tU+\Delta_gU+
		\begin{pmatrix}2&1\\1&2\end{pmatrix}U=0,
		\qquad U(0,x)=\begin{pmatrix}1\\0\end{pmatrix}.
	\]
	空间常值解为 \(U(t,x)=e^{-2t}(\cosh t,-\sinh t)^{\mathsf T}\)。主部强抛物，零阶矩阵甚至正定，但第二分量立即为负。因此一般系统不能逐分量套用标量的次序比较；需要控制几何流中的长度或能量密度时，应先推导该标量量自身满足的微分不等式。
\end{example}
% 纲要条目排为编号小节标题，并进入目录与书签；说明段仍保留于纲要注释中。

% !TeX root = ../../Book3.tex
% 纲要标题：### 26.2 拟线性短时存在
% 纲要位置：几何学计划纲要.md，第 3209 行。
% 正文状态：初稿；尚未正式审核。

% 纲要内容开始
% 
% * 强抛物性及线性化
% * 短时解的构造
% * 短时唯一性和光滑依赖
% * 延拓判据与几何先验界
% 
% 纲要内容结束

\section{拟线性短时存在}
\label{g3:sec:2602}

在线性方程中，系数事先给定；几何流的主系数却常由正在演化的度量或图像决定。例如度量的逆矩阵会出现在二阶导数前，图像方程的主系数还会依赖一阶导数。本节在固定丛上处理这两种情形。先用光滑初值作一个不随时间变化的近似解，利用线性初值同构修正它；短时间的作用，是使近似解的方程误差在时间积分范数中变小。

\ProseParagraphBreak
固定正维闭光滑黎曼流形 \((M^n,\bar g)\) 和正秩实光滑度量丛 \(E\)，所有范数及协变导数使用这些背景数据。整数 \(s>n/2+3\) 固定；用上一节的 \(X_T^s(E)\) 作为未知量空间，用 \(Y_T^{s-1}=L^2(0,T;H^{s-1}(E))\) 作为方程空间。

\subsection{强抛物性及线性化}
\label{g3:subsec:2602:01}
考虑自治方程
\begin{equation}\label{g3:eq:quasilinear-parabolic-system}
	\partial_tu=\mathcal Q(u),\qquad
	\mathcal Q(u)=A^{ij}(x,u,\nabla u)\nabla_i\nabla_j u+B(x,u,\nabla u).
\end{equation}
这里 \(A^{ij}\) 取值于 \(\operatorname{End}(E)\)，\(B\) 取值于 \(E\)，均光滑定义在一阶喷射丛的一个开集 \(\mathcal O\) 上。局部表达在丛坐标变换下组成同一个算子。主系数可假定关于 \(i,j\) 对称，因为反对称部分只通过联络曲率产生零阶项。初值为一个光滑截面 \(u_0\)，其一阶喷射图 \(j^1u_0(M)\) 包含在 \(\mathcal O\) 内。

\begin{definition}\label{g3:def:quasilinear-strong-parabolicity}
	设 \((M^n,\bar g)\) 为正维闭光滑黎曼流形，\(E\to M\) 为带固定度量联络的正秩实光滑度量丛，\(\mathcal O\subset J^1E\) 开，\(A^{ij}\) 与 \(B\) 为其上的光滑系数，分别取值于 \(\operatorname{End}(E)\) 与 \(E\)。记 \(\mathcal Q(u)=A^{ij}(x,u,\nabla u)\nabla_i\nabla_ju+B(x,u,\nabla u)\)。
	若存在 \(\theta>0\)，使一个喷射集合上所有 \(\xi\in\mathrm T_x^*M\)、\(v\in E_x\) 满足
	\[
		\langle A^{ij}(x,z,p)\xi_i\xi_jv,v\rangle
		\ge\theta|\xi|_{\bar g}^2|v|^2.
	\]
	则称式\eqref{g3:eq:quasilinear-parabolic-system}沿该喷射集合一致\term[qiang paowu]{强抛物}[strongly parabolic]。本节的短时条件是此式在 \(j^1u_0(M)\) 上成立。紧性与连续性保证缩小喷射邻域后仍有共同的正下界。
\end{definition}

定义 \(\mathscr P(u)=u_t-\mathcal Q(u)\)。对光滑 \(u\)，其线性化作用于 \(v\) 为
\begin{equation}\label{g3:eq:quasilinear-parabolic-linearization}
	\begin{split}
	\mathrm D\mathscr P(u)v={}&v_t-A^{ij}(j^1u)\nabla_i\nabla_jv\\
	&-\bigl(A_z^{ij}(j^1u)v+A_p^{ij}(j^1u)\nabla v\bigr)\nabla_i\nabla_j u
	-B_z(j^1u)v-B_p(j^1u)\nabla v.
	\end{split}
	\end{equation}
约定 \(j^1u=(x,u,\nabla u)\)。最高阶项仍是 \(v_t-A^{ij}\nabla_i\nabla_jv\)；对一阶喷射求导产生的是关于 \(v\) 的一阶项。因而强抛物性恰给出上一节线性定理所需的主符号条件。

\begin{lemma}\label{g3:lem:quasilinear-parabolic-smooth-map}
	设 \((M^n,\bar g)\) 为正维闭光滑黎曼流形，\(E\to M\) 为带固定度量联络的正秩实光滑度量丛，\(\mathcal O\subset J^1E\) 开，\(A^{ij}\) 与 \(B\) 为其上的光滑系数，分别取值于 \(\operatorname{End}(E)\) 与 \(E\)。记 \(\mathcal Q(u)=A^{ij}(x,u,\nabla u)\nabla_i\nabla_ju+B(x,u,\nabla u)\)。固定整数 \(s>n/2+3\)，\(X_T^s\) 取\cref{g3:def:parabolic-energy-space}的空间，\(Y_T^{s-1}=L^2(0,T;H^{s-1}(E))\)，\(\mathscr P(u)=\partial_tu-\mathcal Q(u)\)。
	在 \(X_T^s(E)\) 中喷射图留在 \(\mathcal O\) 的开集上，\(\mathscr P\) 是到 \(Y_T^{s-1}\) 的光滑映射。对固定光滑截面 \(u_0\) 的一个有界小邻域，各阶微分的界可对 \(0<T\le1\) 一致选取。
\end{lemma}
\begin{proof}
	因为 \(s-1>n/2\)，\(H^{s-1}\) 是乘法代数。由\cref{g3:thm:sobolev-composition-smooth}，映射 \(u\mapsto A(j^1u),B(j^1u)\) 从 \(H^s\) 到 \(H^{s-1}\) 光滑；在连续时间空间上逐点应用同一估计，得到到 \(C_tH^{s-1}\) 的光滑映射。另一方面 \(\nabla^2u\in L^2_tH^{s-1}\)，且
	\[
		\|av\|_{L^2H^{s-1}}
		\le C\|a\|_{C_tH^{s-1}}\|v\|_{L^2H^{s-1}}.
	\]
	这证明主项的光滑性；\(B\) 项从连续时间范数进入 \(L^2\) 范数只产生 \(\sqrt T\le1\)。微分各次使用相同的乘法，余项由积分形式的 Taylor 公式控制，故不是只有形式导数。喷射像的条件由 \(C_tH^s\to C_tC^1\) 保证为开条件。固定 \(u_0\) 的时间常值延拓在 \(X_T^s\) 中有与 \(T\le1\) 无关的界，所以各次复合与乘法常数均可统一。
\end{proof}

\begin{example}\label{g3:exm:graph-parabolic-symbol}
	在平坦环面上，实函数方程
	\[
		u_t=\left(\delta^{ij}-\frac{u_i u_j}{1+|\mathrm Du|^2}\right)u_{ij}
	\]
	具有图像型主部。对任意 \(\xi\)，主二次型介于 \(|\xi|^2/(1+|\mathrm Du|^2)\) 与 \(|\xi|^2\) 之间。因此梯度有界就保证统一的强抛物常数；解的零阶大小在这个符号计算中没有作用。该方程的几何解释将在平均曲率流一章给出。
\end{example}

\subsection{短时解的构造}
\label{g3:subsec:2602:02}
把初值约束先放进未知量空间：写 \(u=u_0+v\)，要求 \(v\in X_{T,0}^s\)。这个选择使所有迭代自动具有同一初值。

\begin{theorem}\label{g3:thm:quasilinear-parabolic-short-time}
	设 \((M^n,\bar g)\) 为正维闭光滑黎曼流形，\(E\to M\) 为带固定度量联络的正秩实光滑度量丛，\(\mathcal O\subset J^1E\) 开，\(A^{ij}\) 与 \(B\) 为其上的光滑系数，分别取值于 \(\operatorname{End}(E)\) 与 \(E\)。记 \(\mathcal Q(u)=A^{ij}(x,u,\nabla u)\nabla_i\nabla_ju+B(x,u,\nabla u)\)。\(u_0\) 光滑且 \(j^1u_0(M)\subset\mathcal O\)，方程在此喷射图上一致强抛物。则存在 \(T>0\) 和唯一的光滑短时解
	\[
		u\in C^\infty([0,T]\times M;E),\qquad
		u_t=\mathcal Q(u),\quad u(0)=u_0.
	\]
	唯一性针对具有该初值的光滑解在共同存在区间内成立。若系数与光滑初值依赖有限维光滑参数，则在每个固定参数附近可选共同的 \(T\)，解也联合光滑依赖这些参数。光滑性包括初时刻的单侧时间导数。
\end{theorem}
\begin{proof}
	先求 \(X_T^s\) 解；其光滑性、唯一性及参数结论在下一小节逐项证明。令 \(F_T(v)=\mathscr P(u_0+v)\)。在 \(v=0\) 处，线性化 \(\mathcal A_T=\mathrm DF_T(0)\) 是具有光滑、与时间无关系数的强抛物算子，作用在零初值空间上。\Cref{g3:thm:linear-parabolic-isomorphism}给出其逆 \(G_T:Y_T^{s-1}\to X_{T,0}^s\)，且 \(\|G_T\|\le C_0\) 对 \(0<T\le1\) 一致。

	由上一引理，存在与 \(T\) 无关的小半径 \(\rho>0\)，使 \(\|v\|_{X_T^s}\le\rho\) 时喷射仍在所选邻域，并且
	\[
		\|G_T(\mathrm DF_T(v)-\mathcal A_T)\|\le\tfrac12.
	\]
	这里先选 \(\rho\) 控制二阶微分界，再缩短 \(T\)，不以“时间短”替代最高阶范数中的小邻域条件。令
	\[
		\mathcal T_T(v)=v-G_TF_T(v).
	\]
	积分形式的微分中值公式给出 \(\mathcal T_T\) 在此球上的 Lipschitz 常数至多为 \(1/2\)。初始误差为
	\[
		\|F_T(0)\|_{Y_T^{s-1}}
		=\sqrt T\,\|\mathcal Q(u_0)\|_{H^{s-1}}.
	\]
	选 \(T\) 使 \(C_0\sqrt T\|\mathcal Q(u_0)\|_{H^{s-1}}\le\rho/2\)，则
	\(\|\mathcal T_T(v)\|\le\rho/2+\|v\|/2\le\rho\)。\Cref{g3:lem:banach-contraction}遂给出唯一不动点，亦即 \(F_T(v)=0\)。可把误差再缩小一半，使解位于球的内部。上述微分估计还说明 \(\mathrm DF_T(v)\) 是 \(\mathcal A_T\) 的可逆小扰动，因此在所得解处仍可使用带参数隐函数定理。
\end{proof}

\begin{example}\label{g3:exm:reaction-diffusion-short-time}
	在正维闭光滑黎曼流形 \((M,\bar g)\) 上，记 \(\Delta_{\bar g}=-\operatorname{div}_{\bar g}\nabla\)。方程 \(u_t+\Delta_{\bar g}u=u^2\) 的主部始终强抛物，所以每个光滑初值都有短时解。但若 \(u_0\equiv a>0\)，解为 \(u(t)=a/(1-at)\)，在 \(t=a^{-1}\) 发散。固定的正主符号保证局部可解，并不限制零阶反应项造成的增长。
\end{example}

\subsection{短时唯一性和光滑依赖}
\label{g3:subsec:2602:03}
先处理参数和正则性。直接在方程中反复求空间导数，会引入高阶系数；也可以先把空间微分改写成对光滑坐标变化参数的微分，使隐函数定理同时完成正则性证明。

\begin{lemma}\label{g3:lem:parabolic-parameter-regularity}
	设 \((M^n,\bar g)\) 为正维闭光滑黎曼流形，\(E\to M\) 为带固定度量联络的正秩实光滑度量丛，\(\mathcal O\subset J^1E\) 开，\(A^{ij}\) 与 \(B\) 为其上的光滑系数，分别取值于 \(\operatorname{End}(E)\) 与 \(E\)。记 \(\mathcal Q(u)=A^{ij}(x,u,\nabla u)\nabla_i\nabla_ju+B(x,u,\nabla u)\)。设整数 \(s>n/2+3\)，\(u_0\) 光滑且其喷射图包含于 \(\mathcal O\)，主符号在该图上一致强抛物。由短时构造得到的 \(X_T^s\) 解在 \([0,T]\times M\) 上光滑。有限维光滑数据族在同一个足够短的时间区间上具有联合光滑解族。
\end{lemma}
\begin{proof}
	对有限维参数 \(a\)，把初值的时间常值延拓记为 \(u_{0,a}\)，写成零初值方程 \(F(a,v)=0\)。光滑数据族在每个固定 Sobolev 范数中光滑，故上一光滑映射引理使 \(F\) 为 Banach 空间间的光滑映射。在已得解处 \(v\) 微分可逆，所以\cref{g3:cor:banach-implicit-parameters}给出 \(v=v(a)\) 在 \(X_T^s\) 中的光滑依赖。

	现在让参数描述源上的坐标变化。取一个光滑向量场 \(Z\) 的流 \(\phi_\varepsilon\)。沿流线用固定联络平行移动，将 \(u(\phi_\varepsilon(x),t)\) 搬回 \(E_x\)，记所得线性作用为 \(U_\varepsilon u\)。换元、链式法则及光滑稠密性说明：它在 \(X_T^s\) 上局部一致有界，且对每个固定 \(u\) 强连续。此处尚不假定它对任意 \(u\) 可微。

	变换后的截面满足
	\[
		w_t=\mathcal Q_\varepsilon(w),\qquad
		\mathcal Q_\varepsilon=U_\varepsilon\mathcal Q U_\varepsilon^{-1},\qquad
		w(0)=U_\varepsilon u_0.
	\]
	直接作坐标变换可见 \(\mathcal Q_\varepsilon\) 仍具有式\eqref{g3:eq:quasilinear-parabolic-system}的形式，其系数关于 \((\varepsilon,x,z,p)\) 光滑；主符号由可逆坐标变换及纤维等距共轭得到，故小 \(\varepsilon\) 时保持统一强抛物。由于 \(u_0\) 光滑，初值 \(U_\varepsilon u_0\) 关于 \(\varepsilon\) 在所有所需范数中光滑。前一段遂构造一个在 \(X_T^s\) 中光滑依赖 \(\varepsilon\) 的解。另一方面，强连续性使 \(U_\varepsilon u\) 留在局部唯一性邻域，故它就是这个解。

	因此 \(\varepsilon\mapsto U_\varepsilon u\) 实际上在 \(X_T^s\) 中光滑。对光滑测试截面作换元后求导，识别其在零点的分布导数为 \(\nabla_Zu\)。所以 \(\nabla_Zu\in X_T^s\)。取有限多个全局光滑向量场，使它们在各坐标小片上包含坐标方向；这样的场可用截断延拓构造。再使用这些流的有限多参数复合，任意阶参数导数就给出任意多次协变导数，外加已经控制的低阶光滑系数项。归纳得到
	\[
		u\in C([0,T];H^{s+k}(E))\qquad(k=0,1,2,\ldots).
	\]
	同一固定时间 \(T\) 适用于所有阶数，因为参数隐函数定理给出的是整个参数邻域内的光滑映射，而不是逐阶重新构造一个解。

	方程于是给出 \(u_t=\mathcal Q(u)\in C_tH^k\) 对每个 \(k\) 成立。反复对方程求时间导数，链式法则依次给出所有混合导数，Sobolev 嵌入把这些结论转为联合光滑性。把数据参数与上述坐标变化参数放在一起，完全相同的论证还证明参数、空间、时间的联合光滑性。
\end{proof}

\begin{proposition}\label{g3:prop:quasilinear-parabolic-uniqueness}
	设 \((M^n,\bar g)\) 为正维闭光滑黎曼流形，\(E\to M\) 为带固定度量联络的正秩实光滑度量丛，\(\mathcal O\subset J^1E\) 开，\(A^{ij}\) 与 \(B\) 为其上的光滑系数，分别取值于 \(\operatorname{End}(E)\) 与 \(E\)。记 \(\mathcal Q(u)=A^{ij}(x,u,\nabla u)\nabla_i\nabla_ju+B(x,u,\nabla u)\)。
	式\eqref{g3:eq:quasilinear-parabolic-system}具有相同初值的两个光滑强抛物解，在共同存在的时间区间上相同。
\end{proposition}
\begin{proof}
	设两个解为 \(u,v\)，\(w=u-v\)。在某个两者相等的时刻后的一小段时间内，它们的一阶喷射位于同一组凸坐标邻域。以 \(A(j^1u)\nabla^2w\) 为主项，用积分中值公式展开系数差，得到
	\[
		w_t=A^{ij}(j^1u)\nabla_i\nabla_jw+C^i(t,x)\nabla_iw+D(t,x)w,
	\]
	其中 \(C,D\) 光滑；例如 \((A(j^1u)-A(j^1v))\nabla^2v\) 只对 \(w\) 含一阶导数。主部直接取自 \(u\)，故保持强抛物，不必假设整条喷射线段上的主部具有同一符号。上一节的零阶能量估计与零初值给出 \(w=0\)。从 \(t=0\) 开始，若相等的时间区间有一个小于共同终点的右端，连续性使两解在该端仍相等，再由同一论证向后延长，矛盾。
\end{proof}

这补齐了短时定理的唯一性与光滑性。对几何未知量，可以先选固定背景下的局部图册。例如映射 \(f\) 接近一个光滑 \(f_0:M\to N\) 时，写 \(f(x)=\exp_{f_0(x)}v(x)\)，未知量 \(v\) 是固定丛 \(f_0^*\mathrm TN\) 的截面。若变换后的方程具有上述强抛物主部，短时解便由本节定理给出；在图册交叠处用唯一性识别所得几何解。

\subsection{延拓判据与几何先验界}
\label{g3:subsec:2602:04}
延拓时必须区分系数的定义域、主符号的正下界和解的导数界。下面给出一个可直接检验的充分判据：控制二阶空间导数，并保证一阶喷射不离开一个强抛物紧区域。首先补出从低阶界提升高阶能量所需的乘积估计。

\begin{lemma}\label{g3:lem:parabolic-tame-products}
	在固定闭流形及有限坐标覆盖上，对整数 \(r\ge1\)，光滑有限分量函数满足
	\begin{align*}
		\|F(\cdot,z)\|_{H^r}&\le C(1+\|z\|_{H^r}),\\
		\|[\nabla^r,a]\nabla^2u\|_{L^2}
		&\le C_r\bigl(\|\nabla a\|_\infty\|u\|_{H^{r+1}}
		+\|a\|_{H^r}\|u\|_{C^2}\bigr).
	\end{align*}
	第一式假定 \(F\) 光滑且 \(z\) 的值在固定紧集内，常数依赖此紧集及 \(F\) 的有限阶导数；第二式的交换子只指 Leibniz 展开中至少一个导数落到 \(a\) 的项。固定联络产生的低阶项可另加 \(C\|a\|_{H^r}\|u\|_{C^2}+C\|a\|_{C^1}\|u\|_{H^{r+1}}\) 吸收。
\end{lemma}
\begin{proof}
	所需插值可以在坐标中直接证明。对紧支光滑 \(f\)，令 \(A_j=\|\mathrm D^jf\|_{L^{2r/j}}\)、\(A_0=\|f\|_\infty\)。在 \(\int|\mathrm D^jf|^{2r/j}\) 中把一个导数分部积分，导数可落到另一高阶因子或其模的幂；Hölder 指数满足
	\[
		\frac{j-1}{2r}+\frac{j+1}{2r}
		+\frac{(2r/j)-2}{2r/j}=1.
	\]
	因而 \(A_j^2\le C_rA_{j-1}A_{j+1}\)（\(1\le j<r\)）。零点处用 \((|\mathrm D^jf|^2+\varepsilon)^{1/2}\) 正则化后取极限。令 \(b_j=\log A_j-(1-j/r)\log A_0-(j/r)\log A_r\)，再从 \(b_j\) 减去屏障 \(\tfrac12(\log C_r)j(r-j)\)。所得数列在内部不超过两个相邻值的平均，端点为零，故离散最大值原理给出
	\[
		\|\mathrm D^jf\|_{L^{2r/j}}
		\le C_r\|f\|_\infty^{1-j/r}\|\mathrm D^rf\|_{L^2}^{j/r}.
	\]
	零端点范数的情形可先加正数再取极限。乘以固定截断并用有限覆盖，在闭流形上将末项换为 \(\|f\|_{H^r}\)，得到同样的局部估计。

	对 \(F(x,z)\) 用链式法则，每项为有界系数乘以若干 \(\mathrm D^{j_i}z\)，且 \(\sum j_i\le r\)。上式和 Hölder 不等式使其 \(L^2\) 范数至多为 \(C(1+\|z\|_{H^r})\)：最高阶总次数为 \(r\) 时高阶范数的幂恰为一，低于 \(r\) 时再用有限体积及 \(a^\eta\le1+a\)、\(0\le\eta\le1\)。这证明第一式。

	第二式的典型项是 \(\mathrm D^\ell a\,\mathrm D^{r-\ell+2}u\)、\(1\le\ell\le r\)。\(r=1\) 直接估计；\(r>1\) 时，把 \(\mathrm Da\) 与 \(\mathrm D^2u\) 分别作为上述插值中的函数，阶数取 \(r-1\)。所用导数阶为 \(\ell-1\) 与 \(r-\ell\)，两者之和正好为 \(r-1\)。Hölder 后再用加权算术几何平均不等式，得到
	\[
		C\bigl(\|\mathrm Da\|_\infty\|\mathrm D^2u\|_{H^{r-1}}
		+\|\mathrm Da\|_{H^{r-1}}\|\mathrm D^2u\|_\infty\bigr).
	\]
	求和并加入固定坐标、联络与截断的低阶项，得到结论。
\end{proof}

\begin{theorem}\label{g3:thm:quasilinear-parabolic-continuation}
	设 \((M^n,\bar g)\) 为正维闭光滑黎曼流形，\(E\to M\) 为带固定度量联络的正秩实光滑度量丛，\(\mathcal O\subset J^1E\) 开，\(A^{ij}\) 与 \(B\) 为其上的光滑系数，分别取值于 \(\operatorname{End}(E)\) 与 \(E\)。记 \(\mathcal Q(u)=A^{ij}(x,u,\nabla u)\nabla_i\nabla_ju+B(x,u,\nabla u)\)。
	设 \(u\) 是式\eqref{g3:eq:quasilinear-parabolic-system}在 \([0,T_*)\) 上的光滑解，\(0<T_*<\infty\)。若存在紧集 \(K\Subset\mathcal O\)，使全部 \(j^1u(x,t)\in K\)，主系数在 \(K\) 上一致强抛物，并且
	\[
		\sup_{0\le t<T_*}\|u(t)\|_{C^2}<\infty,
	\]
	则解可光滑延拓到某个 \([0,T_*+\varepsilon)\)。
\end{theorem}
\begin{proof}
	记给定的 \(C^2\) 界为 \(R\)。先证明每个高阶 Sobolev 范数在 \([0,T_*)\) 上有界。对方程施以 \(\nabla^r\)，主项为 \(A^{ij}(j^1u)\nabla_i\nabla_j(\nabla^ru)\)。其作用在张量丛 \((\mathrm T^*M)^{\otimes r}\otimes E\) 上仍强抛物。上一节零阶 Gårding 估计的冻结系数证明只需要主系数的统一连续模与一阶导数界；这里二者都由 \(K,R\) 控制。因此分部积分给出负的 \(\|\nabla^{r+1}u\|_{L^2}^2\) 主项，余下低阶项受 \(C_r\|u\|_{H^r}^2\) 控制。

	需要特别检查的是导数落到未知系数的项。前一引理给出
	\[
		\|A(j^1u)\|_{H^r}+\|B(j^1u)\|_{H^r}
		\le C_{r,K}(1+\|u\|_{H^{r+1}}),
	\]
	且 \(\|\nabla A(j^1u)\|_\infty\le C_{K,R}\)。故交换子及反应项与 \(\nabla^ru\) 的配对，均至多为
	\(C_{r,K,R}(1+\|u\|_{H^{r+1}})\|u\|_{H^r}\)。Young 不等式吸收一小部分最高阶项。对从零到 \(r\) 阶的能量相加，得到
	\begin{equation}\label{g3:eq:quasilinear-high-energy}
		\frac{\mathrm d}{\mathrm dt}\|u\|_{H^r}^2
		+c_r\|u\|_{H^{r+1}}^2
		\le C_{r,K,R}(1+\|u\|_{H^r}^2).
	\end{equation}
	固定几何曲率的交换项只涉及更低阶导数，已经包含在右端。Gronwall 不等式与光滑初值使每个 \(H^r\) 范数在有限区间内一致有界。

	方程及这些高阶界又给出 \(\sup_{t<T_*}\|u_t(t)\|_{H^r}<\infty\)，因为此处可用已经控制的 \(H^{r+2}\) 范数。于是 \(u(t)\) 在每个 \(H^r\) 中关于时间一致 Lipschitz，故在 \(t\uparrow T_*\) 时收敛到同一个光滑截面 \(u_*\)。喷射极限仍在 \(K\) 内，强抛物性保留。以 \(u_*\) 为初值重新应用短时存在定理，便得到右侧的光滑解。

	在接合时刻，两侧的全部空间导数趋向相同极限；时间导数由方程确定，反复微分方程得到的所有初始时间喷射也一致。因此接合后的解光滑。若未知量取值于流形或开张量空间，先由零阶与时间导数界取得极限，再用有限组局部坐标或极限对象的指数图表示未知量；上述论证在这些固定图中适用，唯一性保证交叠相容。
\end{proof}

该判据也指出几何估计应控制什么。张量方程要保持正定性等开放条件；映射方程要防止像逃出可用坐标区域，并控制空间导数。对调和映射热流，下一节的曲率符号给出一阶界，再由协变导数演化取得所有高阶界，恰能用于上述延拓过程。
% 纲要条目排为编号小节标题，并进入目录与书签；说明段仍保留于纲要注释中。

% !TeX root = ../../Book3.tex
% 纲要标题：### 26.3 调和映射热流的基本演化
% 纲要位置：几何学计划纲要.md，第 3216 行。
% 正文状态：初稿；尚未正式审核。

% 纲要内容开始
% 
% * 张力场热方程及几何不变性
% * 能量下降和时间导数的积分控制
% * 梯度 Bochner 演化
% * 非正目标曲率的符号作用
% 
% 纲要内容结束

\section{调和映射热流的基本演化}
\label{g3:sec:2603}

调和映射的张力场给出能量下降最快的方向。把这个方向作为速度，就得到一个随时间改变映射的方程。短时理论首先提供光滑解；为了继续演化，还要找出哪些量不会失控。本节分别计算总能量、空间微分和时间速度的变化。目标的非正截面曲率将在后两项中给出有利的符号。

固定闭连通光滑黎曼流形 \((M^m,g)\)，\(m\ge1\)，以及紧无边界光滑黎曼流形 \((N^n,h)\)，\(n\ge1\)。源度量 \(g\) 不随时间变化，\(\Delta_g=-\operatorname{div}_g\nabla\)。在解的光滑存在区间 \([0,T)\) 上，记 \(u_i=\mathrm du(e_i)\)、\(V=\partial_tu\)，其中 \((e_i)\) 是与时间无关的源正交标架；拉回联络的时间和空间微分分别记为 \(D_t,D_i\)。

\subsection{张力场热方程及几何不变性}
\label{g3:subsec:2603:01}
第18.1节定义了 \(\tau(u)=\operatorname{tr}_g\nabla\mathrm du\)，并证明能量第一变分为 \(-\int_M\langle\tau(u),V\rangle\mu_g\)。因此负梯度方向正是 \(\tau(u)\)，而不是它的相反数。

\begin{definition}\label{g3:def:harmonic-map-heat-flow}
	设 \((M,g)\) 为闭光滑黎曼流形，\((N,h)\) 为无边界光滑黎曼流形，\(u_0:M\to N\) 光滑。满足
	\begin{equation}\label{g3:eq:harmonic-map-heat-flow}
		\partial_tu=\tau(u),\qquad u(0,x)=u_0(x)
	\end{equation}
	的光滑映射 \(u:[0,T)\times M\to N\) 称为从 \(u_0\) 出发的\term[tiaohe yingshe reliu]{调和映射热流}[harmonic map heat flow]。
\end{definition}

在一张目标坐标图中，式\eqref{g3:eq:tension-coordinates}把方程写成
\begin{equation}\label{g3:eq:harmonic-heat-local-equation}
	\partial_tu^\alpha+\Delta_g u^\alpha
	=g^{ij}\Gamma^\alpha_{\beta\gamma}(h)(u)
	\partial_i u^\beta\partial_j u^\gamma.
\end{equation}
最高阶的空间部分是各分量共同的 \(\Delta_g\)，因而是强抛物的；目标的几何进入一阶导数的二次项。这个坐标表达只在像落入该图时使用，方程本身由张力场内在定义。

\begin{corollary}\label{g3:cor:harmonic-heat-short-time}
	设 \((M,g)\) 为闭光滑黎曼流形，\((N,h)\) 为紧无边界光滑黎曼流形，\(u_0:M\to N\) 光滑。则存在 \(T>0\) 及唯一的光滑调和映射热流 \(u:[0,T)\times M\to N\)，满足 \(u(0)=u_0\)。此结论不要求目标曲率具有特定符号。
\end{corollary}
\begin{proof}
	在固定向量丛 \(E=u_0^*\mathrm TN\) 中取小截面 \(v\)，以
	\[
		F(x,v(x))=\exp_{u_0(x)}v(x)
	\]
	表示邻近 \(u_0\) 的映射。目标紧，故可取统一的正半径，使纤维微分 \(J(x,v)=\mathrm D_vF(x,v)\) 可逆。对 \(u(t,x)=F(x,v(t,x))\)，有 \(\partial_tu=J(x,v)\partial_tv\)。使用 \(E\) 上固定的拉回联络，空间二阶微分的最高阶项为
	\[
		\nabla_i\mathrm du(e_j)=J(x,v)\nabla_i\nabla_jv
		+\text{只依赖于 }x,v,\nabla v\text{ 的光滑项}.
	\]
	这里对 \(F\) 求导两次，只有两次导数都落在 \(v\) 上时才含 \(\nabla^2v\)；其余项至多使用 \(v\) 的一阶导数。取源迹并乘 \(J^{-1}\)，得到固定丛上的方程
	\[
		\partial_tv=g^{ij}\nabla_i\nabla_jv+\mathcal B(x,v,\nabla v),
		\qquad v(0)=0,
	\]
	其中 \(\mathcal B\) 光滑，主部符号为 \(|\xi|_g^2\operatorname{Id}_E\)。\cref{g3:thm:quasilinear-parabolic-short-time}给出光滑短时解。缩短时间使 \(\|v\|_{C^0}\) 留在上述半径内，便得到取值于 \(N\) 的解。任何另一条从 \(u_0\) 出发的光滑解，在足够短的时间内也可写成这样的指数图截面，故由固定丛方程的唯一性与之相同；再从相同时刻重新使用此结论，得到整个共同存在区间上的唯一性。
\end{proof}

\begin{proposition}\label{g3:prop:harmonic-heat-naturality}
	设 \(u:[0,T)\times(M,g)\to(N,h)\) 为光滑调和映射热流。若 \(a:M'\to M\) 为光滑微分同胚，\(b:(N,h)\to(N',h')\) 为等距微分同胚，则 \(b\circ u\circ a\) 是从 \(b\circ u_0\circ a\) 出发、源度量为 \(a^*g\) 的调和映射热流。对常数 \(c>0\)，\(\widetilde u(t,x)=u(t/c,x)\) 是源度量为 \(cg\)、存在于 \([0,cT)\) 的热流。
\end{proposition}
\begin{proof}
	拉回源联络和等距映射推前的目标联络都保持无挠性与度量相容性，故由 Levi--Civita 联络的唯一性，与相应度量的联络相同。因此二阶协变微分按这两个映射变换，取迹得
	\[
		\tau_{a^*g,h'}(b\circ u\circ a)
		=\mathrm db\bigl(\tau_{g,h}(u)\bigr)\circ a.
	\]
	右端也是 \(b\circ u\circ a\) 的时间导数。常数缩放不改变源联络，但逆度量变为 \(c^{-1}g^{-1}\)，所以 \(\tau_{cg,h}(u)=c^{-1}\tau_{g,h}(u)\)，与 \(u(t/c,x)\) 的时间导数一致。
\end{proof}

时间无关是上述重参数化的条件。若 \(a_t\) 随时间变化，\(u(t,a_t(x))\) 的速度还包含 \(\mathrm du(\partial_ta_t)\)，一般不再满足原来的方程。这也说明调和映射热流对固定源度量演化时，不能任意重新选择随时间变化的源参数。

\begin{example}\label{g3:exm:circle-map-heat-flow}
	令源和目标都是 \(\mathbb R/(2\pi\mathbb Z)\)，取标准度量。对整数 \(d\) 和实数 \(a,b\)，
	\[
		u(t,x)=[dx+b+a e^{-t}\sin x]
	\]
	是一条光滑热流。方括号表示模 \(2\pi\) 的类：\(x\) 增加 \(2\pi\) 时，括号内增加 \(2\pi d\)，所以映射良定义。在局部提升中，目标联络为零，方程就是 \(u_t=u_{xx}\)。周期振荡衰减，次数为 \(d\) 的线性部分保持不变。这条流在任意阶都收敛到调和映射 \([dx+b]\)。
\end{example}

\subsection{能量下降和时间导数的积分控制}
\label{g3:subsec:2603:02}
总能量下降同时控制了速度在全部时间上的平方积分。这个控制不依赖映射是否已经接近某个调和解；它由方程和闭源上的分部积分直接给出。

\begin{theorem}\label{g3:thm:harmonic-heat-energy}
	设 \((M,g)\) 为闭光滑黎曼流形，\((N,h)\) 为无边界光滑黎曼流形，\(u:[0,T)\times M\to N\) 为光滑调和映射热流。定义
	\[
		E(t)=\frac12\int_M|\mathrm du(t)|^2\,\mu_g,
		\qquad Q(t)=\int_M|\partial_tu(t)|^2\,\mu_g.
	\]
	则对 \(0\le a<b<T\)，
	\begin{equation}\label{g3:eq:harmonic-heat-energy-dissipation}
		E'(t)=-Q(t),\qquad E(b)+\int_a^bQ(t)\,\mathrm dt=E(a).
	\end{equation}
	特别地，\(E(t)\) 非增，极限 \(E_* =\lim_{t\uparrow T}E(t)\ge0\) 存在，并且
	\[
		\int_0^T\!\int_M|\partial_tu|^2\,\mu_g\,\mathrm dt
		=E(0)-E_*\le E(0).
	\]
	这里允许 \(T=\infty\)。
\end{theorem}
\begin{proof}
	由\cref{g3:prop:map-covariant-identities}的时间交换式，\(D_tu_i=D_iV\)。源度量固定，故
	\[
		\partial_t\frac12|\mathrm du|^2=\sum_i\langle D_iV,u_i\rangle.
	\]
	在闭源上积分，按照\cref{g3:thm:harmonic-map-first-variation}的散度计算，右端积分等于 \(-\int_M\langle V,\tau(u)\rangle\mu_g\)。代入 \(V=\tau(u)\) 得第一式；按时间积分得到第二式。能量非负且非增，因而存在所述极限；非负积分的单调收敛给出全时间恒等式。
\end{proof}

在\cref{g3:exm:circle-map-heat-flow}中，直接计算可得
\[
	E(t)=\pi d^2+\frac\pi2a^2e^{-2t},\qquad Q(t)=\pi a^2e^{-2t}.
\]
能量损失恰好等于周期振荡的能量，次数所确定的部分没有下降。一般目标上未必有这样的显式分解，但式\eqref{g3:eq:harmonic-heat-energy-dissipation}始终把能量损失与速度联系起来。

若某个时间区间上的能量恒定，则其上 \(Q=0\)，由光滑性得到 \(\partial_tu=0\)。所以能量恒定的光滑热流正是驻定的调和映射。能量只是趋于一个极限时，时间速度的平方可积；将这项积分信息转成点态控制，还需要下面的演化方程。

\subsection{梯度 Bochner 演化}
\label{g3:subsec:2603:03}
能量密度 \(e=|\mathrm du|^2/2\) 是一个标量函数，可以应用最大值原理。它的演化包含源和目标两种曲率：源曲率来自对空间协变导数换序，目标曲率来自拉回丛的联络。

\begin{proposition}\label{g3:prop:harmonic-heat-bochner}
	设 \(u:[0,T)\times(M^m,g)\to(N,h)\) 为光滑调和映射热流，两流形均光滑且无边界。令 \(e=|\mathrm du|^2/2\)，\((e_i)\) 为源的局部正交标架，\(u_i=\mathrm du(e_i)\)。则
	\begin{align}
		(\partial_t+\Delta_g)e
		={}&-|\nabla\mathrm du|^2
		-\sum_i\langle\mathrm du(\operatorname{Ric}_g^\sharp e_i),u_i\rangle\notag\\
		&+\sum_{i,j}\langle R^N(u_i,u_j)u_j,u_i\rangle.
		\label{g3:eq:harmonic-heat-bochner}
	\end{align}
	若 \(\operatorname{Ric}_g\ge-Kg\)、\(K\ge0\)，且 \(N\) 的截面曲率非正，则
	\begin{equation}\label{g3:eq:harmonic-heat-energy-subsolution}
		(\partial_t+\Delta_g)e\le-|\nabla\mathrm du|^2+2Ke.
	\end{equation}
\end{proposition}
\begin{proof}
	空间 Bochner 恒等式\cref{g3:prop:harmonic-map-bochner}在 \(L_g=-\Delta_g\) 的记号下为
	\[
		-\Delta_ge=|\nabla\mathrm du|^2
		+\sum_i\langle D_i\tau(u),u_i\rangle
		+\sum_i\langle\mathrm du(\operatorname{Ric}_g^\sharp e_i),u_i\rangle
		-\sum_{i,j}\langle R^N(u_i,u_j)u_j,u_i\rangle.
	\]
	上一小节已算出 \(\partial_te=\sum_i\langle D_iV,u_i\rangle\)。代入 \(V=\tau(u)\)，两项相消，得到式\eqref{g3:eq:harmonic-heat-bochner}。

	在一点将源 Ricci 张量对角化。每个特征值至少为 \(-K\)，所以源曲率和至少为 \(-K|\mathrm du|^2=-2Ke\)。按照本书约定，
	\[
		\langle R^N(A,B)B,A\rangle
		=\operatorname{Sect}_N(A\wedge B)
		\bigl(|A|^2|B|^2-\langle A,B\rangle^2\bigr)
	\]
	对线性无关的 \(A,B\) 成立，线性相关时左侧为零。非正截面曲率使式\eqref{g3:eq:harmonic-heat-bochner}中每个目标曲率项非正，得到所列不等式。
\end{proof}

粗略地舍去负的导数项，\cref{g3:cor:parabolic-gronwall}给出
\[
	\sup_M e(t)\le e^{2Kt}\sup_M e(0).
\]
它已经排除了任意预先固定的有限时间区间内梯度无界的可能；当 \(K>0\) 时，这个表达式还随时间增长。下一节将把能量积分界与固定时间的热传播结合，得到不随终止时间增长的上界。

\subsection{非正目标曲率的符号作用}
\label{g3:subsec:2603:04}
空间微分的演化中保留了源 Ricci 曲率，时间速度的演化则没有这一项。后者是沿同一源点的目标向量场，空间上不携带额外的源余切指标。这个差别使非正目标曲率直接控制最大速度。

\begin{proposition}\label{g3:prop:harmonic-heat-speed}
	设 \((M,g)\) 为闭光滑黎曼流形，\((N,h)\) 为无边界光滑黎曼流形，\(u:[0,T)\times M\to N\) 为光滑调和映射热流，\(V=\partial_tu\)，\(Q(t)=\int_M|V|^2\mu_g\)。记 \(u_0=u(0,\cdot)\)，\((e_i)\) 为源的局部正交标架，\(u_i=\mathrm du(e_i)\)，\(\tau(u_0)=\operatorname{tr}_g\nabla\mathrm du_0\)。令 \(D^*D\) 是 \(u(t)^*\mathrm TN\) 上的非负联络 Laplace 算子，则
	\begin{equation}\label{g3:eq:harmonic-heat-jacobi-velocity}
		D_tV+D^*DV=\sum_iR^N(V,u_i)u_i.
	\end{equation}
	若 \(\operatorname{Sect}_N\le0\)，则
	\begin{equation}\label{g3:eq:harmonic-heat-speed-subsolution}
		(\partial_t+\Delta_g)|V|^2
		=-2|DV|^2+2\sum_i\langle R^N(V,u_i)u_i,V\rangle
		\le-2|DV|^2.
	\end{equation}
	因而
	\[
		\sup_M|\partial_tu(t)|\le\sup_M|\tau(u_0)|,
		\qquad Q'(t)\le-2\int_M|DV|^2\,\mu_g\le0.
	\]
\end{proposition}
\begin{proof}
	将源标架按与时间无关的方式延拓。由式\eqref{g3:eq:map-variation-commutation}，
	\[
		D_tD_i u_i=D_iD_tu_i+R^N(V,u_i)u_i
		=D_iD_iV+R^N(V,u_i)u_i.
	\]
	在张力场中还要减去 \(\mathrm du(\nabla^M_{e_i}e_i)\)；它的时间导数是 \(D_{\nabla^M_{e_i}e_i}V\)。对 \(i\) 求和，便得
	\[
		D_t\tau(u)=-D^*DV+\sum_iR^N(V,u_i)u_i.
	\]
	再用 \(\tau(u)=V\)，得到式\eqref{g3:eq:harmonic-heat-jacobi-velocity}。

	联络的度量相容性给出
	\[
		\partial_t\frac{|V|^2}{2}=\langle D_tV,V\rangle,
		\qquad
		\Delta_g\frac{|V|^2}{2}=\langle D^*DV,V\rangle-|DV|^2.
	\]
	相加并使用非正截面曲率的符号，得到速度不等式。最大值原理给出 \(\sup|V|^2\) 非增，初速度是 \(\tau(u_0)\)。在闭源上积分，\(\int_M\Delta_g|V|^2\mu_g=0\)，遂得 \(Q'\) 的估计。
\end{proof}

这里空间梯度与时间速度受控的原因不同：前者需要源 Ricci 下界与总能量，后者的最大值仅利用目标曲率符号。若目标曲率有正的部分，式\eqref{g3:eq:harmonic-heat-bochner}会出现正的四次梯度项，式\eqref{g3:eq:harmonic-heat-speed-subsolution}也不再给出无条件的速度单调性。非正曲率将在长期存在证明中同时保留这两种控制。
% 纲要条目排为编号小节标题，并进入目录与书签；说明段仍保留于纲要注释中。

% !TeX root = ../../Book3.tex
% 纲要标题：### 26.4 长期存在的先验控制
% 纲要位置：几何学计划纲要.md，第 3223 行。
% 正文状态：初稿；尚未正式审核。

% 纲要内容开始
% 
% * 闭源上的最大值估计
% * 紧目标下的局部坐标控制
% * 高阶抛物估计
% * 由延拓判据排除有限时奇点
% 
% 纲要内容结束

\section{长期存在的先验控制}
\label{g3:sec:2604}

能量不增给出空间积分控制，Bochner 演化给出逐点微分不等式。把两者通过源流形上的热核联系起来，便能得到与经过时间无关的一阶界。随后对空间导数作带时间权重的估计，将一阶界提升为全部高阶界，最后应用短时理论越过任何有限的时间端点。

本节始终设 \((M^m,g)\) 为闭连通光滑黎曼流形，\((N^n,h)\) 为紧无边界光滑黎曼流形，\(m,n\ge1\)，且 \(\operatorname{Sect}_N\le0\)。令 \(u:[0,T)\times M\to N\) 是光滑初值 \(u_0\) 的光滑调和映射热流。取 \(K\ge0\)，使 \(\operatorname{Ric}_M\ge-Kg\)，并记 \(P=\partial_t+\Delta_g\)。估计先在任意已有的光滑存在区间上建立，因此不会预先假定长期存在。

\subsection{闭源上的最大值估计}
\label{g3:subsec:2604:01}
由最大值原理直接得到的 \(\sup_M e(t)\le e^{2Kt}\sup_M e(0)\) 足以控制有限区间，却不能给出 \(t\to\infty\) 时的统一界。改从稍早的时刻重新比较，可以使用能量的 \(L^1\) 界。

\begin{proposition}\label{g3:prop:harmonic-heat-uniform-gradient}
	设 \((M,g)\) 闭连通，\((N,h)\) 紧、无边界且截面曲率非正，\(u:[0,T)\times M\to N\) 是光滑调和映射热流，\(\operatorname{Ric}_M\ge-Kg\)，\(K\ge0\)。记 \(T_s=e^{-s\Delta_g}\)，并以 \(K_M(s,x,y)\) 表示源流形的标量热核。则对 \(0\le a<t<T\)，有
	\begin{equation}\label{g3:eq:harmonic-heat-kernel-comparison}
		e(t,x)\le e^{2K(t-a)}T_{t-a}(e(a,\cdot))(x).
	\end{equation}
	特别地，令 \(H_1=\max_{M\times M}K_M(1,x,y)\)，则
	\begin{equation}\label{g3:eq:harmonic-heat-gradient-uniform}
		\sup_{0\le t<T}\|\mathrm du(t)\|_{L^\infty}^2
		\le 2e^{2K}\max\{\|e(u_0)\|_{L^\infty},\ H_1E(u_0)\}.
	\end{equation}
	这里右端不依赖 \(T\)。
\end{proposition}

\begin{proof}
	固定 \(a\)。由式\eqref{g3:eq:harmonic-heat-energy-subsolution}，函数 \(w(t)=e^{-2K(t-a)}e(t)\) 满足 \(Pw\le0\)。另一方面，\(z(t)=T_{t-a}e(a)\) 满足 \(Pz=0\)，且 \(w(a)=z(a)\)。\cref{g3:cor:parabolic-gronwall}中的比较原理给出 \(w\le z\)，即式\eqref{g3:eq:harmonic-heat-kernel-comparison}。

	第16.1节证明热核在正时间光滑、非负。因此 \(H_1\) 有限，且对非负函数 \(f\) 有
	\[
		0\le T_1f(x)=\int_MK_M(1,x,y)f(y)\,\mu_g(y)
		\le H_1\|f\|_{L^1}.
	\]
	当 \(1\le t<T\) 时，取 \(a=t-1\)，利用 \(\int_M e(a)\,\mu_g=E(u(a))\le E(u_0)\)，得到
	\[
		e(t,x)\le e^{2K}H_1E(u_0).
	\]
	当 \(0\le t<\min\{1,T\}\) 时，直接的最大值估计给出 \(e(t,x)\le e^{2K}\|e(u_0)\|_{L^\infty}\)。由于 \(|\mathrm du|^2=2e\)，合并两段便得结论。
\end{proof}

这里的时间长度 \(1\) 没有特殊意义。任何固定 \(\theta>0\) 都给出
\[
	\|e(t)\|_{L^\infty}\le e^{2K\theta}
	\max_{x,y}K_M(\theta,x,y)\,E(u_0),\qquad t\ge\theta.
\]
所需事实只是固定正时间的热算子将 \(L^1\) 送入 \(L^\infty\)。这种重新从 \(t-\theta\) 出发的比较避免了把 \(e^{2Kt}\) 一直累积到很大的时间。

\begin{figure}[htbp]
	\centering
	% !TeX root = ../Book3.tex
\begin{tikzpicture}[x=1cm,y=1cm,>=stealth]
	\fill[blue!8] (3.4,-.12) rectangle (6.4,.7);
	\draw[->] (0,0)--(7.3,0) node[right] {时间};
	\draw (.3,-.07)--(.3,.07) node[below=5pt] {\(0\)};
	\draw (3.4,-.07)--(3.4,.07) node[below=5pt] {\(t-1\)};
	\draw (6.4,-.07)--(6.4,.07) node[below=5pt] {\(t\)};
	\draw[<->,thick,blue!65!black] (3.4,.58)--(6.4,.58)
		node[midway,above] {固定长度 \(1\)};
	\node[align=center] at (3.4,-1.05) {\(\|e(t-1)\|_{L^1}\le E(u_0)\)};
	\node[align=center] at (6.4,-1.8) {\(\|e(t)\|_\infty\le e^{2K}H_1E(u_0)\)};
	\draw[->,thick] (3.4,-1.35)--(3.4,-1.8)--(4.15,-1.8);
	\node[anchor=west] at (.3,1.35) {\(e(t)\le e^{2K}T_1e(t-1),\qquad T_1:L^1\longrightarrow L^\infty\)};
\end{tikzpicture}

	\caption[固定时间窗内的热核比较]{每次只在长度为一的时间窗内传播能量密度。窗左端的积分始终受初始能量控制，因此右端的逐点上界不随 \(t\) 增长。图中 \(H_1=\max_{x,y}K_M(1,x,y)\)。}
	\label{g3:fig:harmonic-heat-fixed-window}
\end{figure}

以下用 \(L\) 表示式\eqref{g3:eq:harmonic-heat-gradient-uniform}右端的平方根。再由\cref{g3:prop:harmonic-heat-speed}，记
\begin{equation}\label{g3:eq:harmonic-heat-first-space-time-bounds}
	|\mathrm du|\le L,\qquad |\partial_tu|\le W,
	\qquad W=\|\tau(u_0)\|_{L^\infty}.
\end{equation}
两个常数都只由固定几何与初值决定。

\subsection{紧目标下的局部坐标控制}
\label{g3:subsec:2604:02}
映射值落在紧目标中，并不表示整个像都包含在一张坐标图里。需要控制的是每个小的时空柱中所用坐标的大小和系数界。

\begin{lemma}\label{g3:lem:harmonic-heat-uniform-charts}
	设 \((M,g)\)、\((N,h)\) 为闭光滑黎曼流形，\(u:[0,T)\times M\to N\) 光滑，并满足 \(|\mathrm du|\le L\)、\(|\partial_tu|\le W\)。存在只依赖两流形、\(L,W\) 的 \(r>0\)，使每个柱
	\[
		B_{2r}(x_0)\times[\max\{0,t_0-4r^2\},t_0],\qquad 0\le t_0<T,
	\]
	的像都落在以 \(u(x_0,t_0)\) 为中心的同一大小的目标正规坐标球内。在相应源、目标坐标中，度量、逆度量及任意给定阶的系数导数均有与 \((x_0,t_0)\) 无关的界；映射坐标 \(v\) 的 \(|v|,|\partial_xv|,|\partial_tv|\) 也一致有界。
\end{lemma}

\begin{proof}
	沿最短源测地线，再沿时间方向积分，得
	\begin{equation}\label{g3:eq:harmonic-heat-space-time-distance}
		d_N(u(x,t),u(x_0,t_0))
		\le Ld_M(x,x_0)+W|t-t_0|.
	\end{equation}
	由\cref{g3:lem:uniform-normal-neighborhood}及紧性，可取统一的目标正规坐标半径 \(\rho>0\)，并使半径 \(\rho/2\) 的闭球仍在坐标定义域内。取小于源正规坐标半径一半的 \(r>0\)，并要求
	\[
		2Lr+4Wr^2<\rho/2.
	\]
	式\eqref{g3:eq:harmonic-heat-space-time-distance}便保证整个柱的像落在所需目标球内。

	正规坐标由底点及一个正交标架确定。指数映射对这些参数光滑，紧流形的正交标架丛也紧。因此在固定较小半径上，坐标度量、逆度量及任意固定阶导数统一有界；这与\cref{g3:prop:normal-expansion-uniform-remainders}中的一致性论证相同。坐标映射的一阶导数有界，再结合 \(|\mathrm du|\le L\)、\(|\partial_tu|\le W\)，即得 \(v\) 的三个界。
\end{proof}

在这样的柱内，热流方程是
\[
	\partial_tv^\alpha-g^{ij}\partial_i\partial_jv^\alpha
	=-g^{ij}\Gamma^k_{ij}(g)\partial_kv^\alpha
	+g^{ij}\Gamma^\alpha_{\beta\gamma}(h)(v)
	\partial_iv^\beta\partial_jv^\gamma.
\]
主部由固定源度量给出，统一强抛物；右端的系数在这些图中具有统一界。柱的位置和目标中心可以改变，上述常数却不改变。下一小节在内在张量上估计高阶导数，然后用这些坐标把张量界转成普通导数界。

\subsection{高阶抛物估计}
\label{g3:subsec:2604:03}
一阶界之所以足够，是因为热流的协变微分方程具有受控的最高阶项。先把这种结构写清楚。置
\[
	A_j=\nabla^j\mathrm du,\qquad F_j=|A_j|^2,\qquad j=0,1,\ldots,
\]
其中 \(A_j\) 取值于 \((\mathrm T^*M)^{\otimes(j+1)}\otimes u^*\mathrm TN\)，联络同时使用源 Levi--Civita 联络和目标拉回联络。

\begin{lemma}\label{g3:lem:harmonic-heat-derivative-evolution}
	设 \((M,g)\)、\((N,h)\) 闭且光滑，\(u\) 为光滑调和映射热流。给定整数 \(j\ge1\)，若 \(|A_0|,\ldots,|A_{j-1}|\) 在某个时间区间上一致有界，则存在由这些界和固定几何决定的 \(C_j\)，使该区间上
	\begin{equation}\label{g3:eq:harmonic-heat-higher-evolution}
		PF_j\le-2F_{j+1}+C_j(1+F_j).
	\end{equation}
	若另有 \(\operatorname{Sect}_N\le0\)、\(\operatorname{Ric}_M\ge-Kg\)，则最低阶满足
	\begin{equation}\label{g3:eq:harmonic-heat-base-derivative-evolution}
		PF_0\le-2F_1+2KF_0.
	\end{equation}
\end{lemma}

\begin{proof}
	第二式就是式\eqref{g3:eq:harmonic-heat-energy-subsolution}的两倍。为证明第一式，先考察张量方程。第18.3节的微分交换计算给出
	\begin{equation}\label{g3:eq:harmonic-heat-one-form-evolution}
		(D_t+\nabla^*\nabla)\mathrm du(e_i)
		=\sum_aR^N(u_i,u_a)u_a
		-\mathrm du(\operatorname{Ric}_M^\sharp e_i).
	\end{equation}
	确实，空间二阶迹为
	\[
		\operatorname{tr}_g\nabla^2\mathrm du(e_i)
		=\nabla_i\tau(u)+\sum_aR^N(u_a,u_i)u_a
		+\mathrm du(\operatorname{Ric}_M^\sharp e_i),
	\]
	而 \(D_t\mathrm du(e_i)=\nabla_i\tau(u)\)、\(\nabla^*\nabla=-\operatorname{tr}_g\nabla^2\)，两式相减即得\eqref{g3:eq:harmonic-heat-one-form-evolution}。

	说明反复求导会产生哪些项。以下用星号表示由度量完成的有限次张量积与缩并；它只隐藏排列和整数系数，每一项的范数都不超过各因子范数的乘积乘一个由阶数决定的常数。对有目标向量指标及若干源协变指标的张量 \(A\)，空间交换子由
	\[
		R^N(\mathrm du,\mathrm du)*A\quad\text{及}\quad R^M*A
	\]
	组成；源曲率分别作用在各个源指标上。时间交换子为
	\[
		[D_t,\nabla_i]A=R^N(\tau(u),u_i)A,
	\]
	其中曲率只作用在目标指标上。把空间交换子再求一次导数，可知 \([\nabla^*\nabla,\nabla]A\) 中的项具有以下形式：
	\begin{align*}
		&R^M*\nabla A,\quad \nabla R^M*A,
		\quad R^N*A_0*A_0*\nabla A,\\
		&(\nabla^NR^N)*A_0*A_0*A_0*A,
		\quad R^N*A_1*A_0*A.
	\end{align*}
	这里列出的是项的类型，因而无需借星号规定曲率项的符号。

	令 \(\mathcal R_j=(D_t+\nabla^*\nabla)A_j\)。由式\eqref{g3:eq:harmonic-heat-one-form-evolution}起步，归纳可知 \(\mathcal R_j\) 的每一项都是源曲率或其导数、目标曲率或其协变导数与若干 \(A_{r_1},\ldots,A_{r_q}\) 的缩并，且
	\begin{equation}\label{g3:eq:harmonic-heat-derivative-count}
		r_1+\cdots+r_q\le j.
	\end{equation}
	这个导数计数可直接核验。\(j=0\) 时右端只含 \(A_0\)。对 \(\mathcal R_{j-1}\) 求一次空间导数，若导数落在 \(A_r\) 上，便把 \(r\) 增加 \(1\)；若落在目标曲率上，则增加一个 \(A_0\) 因子；若落在源曲率上，则不增加任何 \(A_r\) 的指标。新增的交换项作用于 \(A_{j-1}\)。其中时间项为
	\[
		R^N*(\operatorname{tr}_gA_1)*A_0*A_{j-1},
	\]
	指标之和恰为 \(j\)；上述空间交换项的指标之和分别至多为 \(j\)。这证明了归纳步骤。

	尤其在 \(j=1\) 时，所有含 \(A_1\) 的反应项都对 \(A_1\) 线性。再求一次导数时，时间交换项中出现的是
	\[
		R^N*A_2*A_0*A_0,\qquad
		R^N*A_1*A_1*A_0,\qquad
		(\nabla^NR^N)*A_1*A_0*A_0*A_0,
	\]
	以及这些类型的排列。\(A_2\) 只线性出现，\(A_1\) 的平方则由上一阶界控制。这也具体说明为何时间联络的变化没有引入更高的导数损失。

	对于一般 \(j\ge1\)，式\eqref{g3:eq:harmonic-heat-derivative-count}保证：若某项包含 \(A_j\)，则该项只包含一个 \(A_j\)，其余 \(A_r\) 都是 \(A_0\)。没有 \(A_j\) 的项由假设中的低阶界控制。两流形紧，曲率各阶导数有界，因此
	\[
		|\mathcal R_j|\le C_j(1+|A_j|).
	\]
	联络保持度量，源度量又不随时间改变，故
	\[
		PF_j=2\langle\mathcal R_j,A_j\rangle-2|\nabla A_j|^2
		\le-2F_{j+1}+C_j(1+F_j),
	\]
	最后用 \(2|A_j|\le1+|A_j|^2\) 并调整常数。这就完成了证明。
\end{proof}

\begin{theorem}\label{g3:thm:harmonic-heat-uniform-derivatives}
	设 \((M,g)\) 闭连通，\((N,h)\) 紧、无边界且截面曲率非正，\(u:[0,T)\times M\to N\) 是光滑初值 \(u_0\) 的光滑调和映射热流。对每个 \(\varepsilon>0\) 及整数 \(j\ge0\)，存在只依赖固定几何、\(u_0,\varepsilon,j\) 的常数 \(C_{j,\varepsilon}\)，使
	\begin{equation}\label{g3:eq:harmonic-heat-uniform-spatial-derivatives}
		\sup_{\varepsilon\le t<T}\|\nabla^j\mathrm du(t)\|_{L^\infty}
		\le C_{j,\varepsilon}.
	\end{equation}
	在\cref{g3:lem:harmonic-heat-uniform-charts}给出的各个固定坐标柱中，以 \(v\) 表示映射的目标坐标。对任意整数 \(a,b\ge0\)，在柱内满足 \(t\ge\varepsilon\) 的部分还有
	\begin{equation}\label{g3:eq:harmonic-heat-uniform-space-time-derivatives}
		\bigl|\partial_t^a\partial_x^\beta v\bigr|
		\le C_{a,b,\varepsilon}.
	\end{equation}
	这里 \(|\beta|=b\)，所有常数均不依赖柱的位置及 \(T\)。
\end{theorem}

\begin{proof}
	\cref{g3:prop:harmonic-heat-uniform-gradient}已给出 \(F_0\) 在整个存在区间的界。固定所要估计的阶数 \(j\ge1\)，取
	\[
		\delta=\min\{1,\varepsilon/(j+1)\}.
	\]
	逐次证明 \(F_p\) 在 \(t\ge p\delta\) 时一致有界，\(p=1,\ldots,j\)。假设已处理所有低于 \(p\) 的阶数。任取 \(p\delta\le t<T\)，置 \(a=t-\delta\)。在 \([a,t]\) 上低阶界均已成立，因此\cref{g3:lem:harmonic-heat-derivative-evolution}给出
	\[
		PF_p\le-2F_{p+1}+C(1+F_p),\qquad
		PF_{p-1}\le-2F_p+C'.
	\]
	第二式在 \(p=1\) 时使用式\eqref{g3:eq:harmonic-heat-base-derivative-evolution}，在 \(p>1\) 时把已经有界的 \(F_{p-1}\) 吸收进 \(C'\)。常数与所选的 \(t\) 无关。

	令 \(s\in[a,t]\)，考虑带时间权重的函数
	\[
		Z(s,x)=(s-a)F_p(s,x)+A F_{p-1}(s,x),
		\qquad A=1+C\delta.
	\]
	于是
	\[
		PZ\le(1+C\delta-2A)F_p+C\delta+AC'
		\le C\delta+AC'.
	\]
	在 \(s=a\) 时，\(Z(a)=AF_{p-1}(a)\) 已有统一界，最高阶项的系数正好为零。\cref{g3:thm:parabolic-maximum}给出
	\[
		\delta F_p(t,x)\le Z(t,x)
		\le A\sup_MF_{p-1}(a)+\delta(C\delta+AC').
	\]
	右端只涉及已经控制的量。这证明归纳步骤。因为 \(j\delta<\varepsilon\)，便得到式\eqref{g3:eq:harmonic-heat-uniform-spatial-derivatives}。这个论证同时说明正时间的高阶控制怎样从低阶控制产生，而不需要先假定任意长区间上的高阶范数有界。

	将内在界换成坐标导数界。在上述柱的固定源、目标坐标中，先有
	\[
		\partial_i\partial_kv^\alpha
		=(\nabla\mathrm du)_{ik}^\alpha
		+\Gamma^\ell_{ik}(g)\partial_\ell v^\alpha
		-\Gamma^\alpha_{\beta\gamma}(h)(v)
		\partial_iv^\beta\partial_kv^\gamma.
	\]
	度量与逆度量的统一界把张量范数界换成分量界，上式于是控制二阶坐标导数。继续求空间导数，每次只增加一个映射导数，或对联络系数求导；归纳得到所有坐标导数的界。各图中联络系数的任意固定阶导数均一致有界，故所得常数与柱的位置无关。这证明式\eqref{g3:eq:harmonic-heat-uniform-space-time-derivatives}的 \(a=0\) 情形。

	坐标中的方程可简写为
	\begin{equation}\label{g3:eq:harmonic-heat-coordinate-time-recursion}
		\partial_tv=g^{ij}\partial_i\partial_jv+Q(x,v,\partial_xv),
	\end{equation}
	其中 \(Q\) 由前一小节的两个一阶项组成，对 \(\partial_xv\) 至多二次，系数光滑且具有所需各阶统一界。该式先控制 \(\partial_t\partial_x^\beta v\)。再对时间求导，每出现一次 \(\partial_tv\) 或其空间导数，就用同一方程代入。归纳得到 \(\partial_t^a\partial_x^\beta v\) 是至多 \(2a+b\) 阶空间导数和固定光滑系数的有限次乘积之和。前面已经控制所有空间阶数，故得到任意混合阶数的界。
\end{proof}

上述坐标计算也适用于任何像始终留在同一较小目标正规坐标球内的固定源坐标邻域。

\subsection{由延拓判据排除有限时奇点}
\label{g3:subsec:2604:04}
\begin{theorem}\label{g3:thm:harmonic-heat-global-existence}
	设 \((M^m,g)\) 为闭连通光滑黎曼流形，\((N^n,h)\) 为紧无边界光滑黎曼流形，\(m,n\ge1\)，且 \(\operatorname{Sect}_N\le0\)。则每个光滑映射 \(u_0:M\to N\) 都确定唯一的光滑调和映射热流
	\[
		u:[0,\infty)\times M\longrightarrow N,\qquad u(0)=u_0.
	\]
	其能量满足\cref{g3:thm:harmonic-heat-energy}的耗散恒等式，一阶导数满足式\eqref{g3:eq:harmonic-heat-first-space-time-bounds}，每个正时间之后的全部高阶导数满足\cref{g3:thm:harmonic-heat-uniform-derivatives}的统一界。
\end{theorem}

\begin{proof}
	\cref{g3:cor:harmonic-heat-short-time}把从 \(u_0\) 出发的解延伸到一个最大光滑存在区间 \([0,T_{\max})\)。短时唯一性保证这些延伸彼此相容。设 \(T_{\max}<\infty\)。首先，由速度界 \(|\partial_tu|\le W\)，有
	\[
		\sup_{x\in M}d_N(u(t,x),u(s,x))\le W|t-s|.
	\]
	因此 \(u(t)\) 在完备目标中一致收敛到连续映射 \(u_*:M\to N\)。围绕每个 \(x\in M\) 选源坐标邻域 \(U_x\)，使 \(u_*(\overline{U_x})\) 落在某个目标正规坐标球的较小闭球内。由一致收敛，充分接近 \(T_{\max}\) 时，\(u(t,\overline{U_x})\) 都留在同一个目标图中。紧性允许只选有限个 \(U_x\)，并选一个共同的起始时刻 \(a<T_{\max}\)。

	在这些固定图中，取 \(\varepsilon=T_{\max}/2\)，由前一定理的内在导数界及其坐标转换得到全部空间导数界；式\eqref{g3:eq:harmonic-heat-coordinate-time-recursion}又给出全部混合导数界。若 \(v\) 表示映射坐标，则对每个整数 \(k\ge0\)，在较小的源坐标闭包上有
	\[
		\|v(t)-v(s)\|_{C^k}\le C_k|t-s|,
		\qquad \max\{a,T_{\max}/2\}\le s,t<T_{\max}.
	\]
	这是对 \(\partial_t\partial_x^\beta v\) 积分、再对 \(|\beta|\le k\) 取上确界的结果。所以坐标函数及各阶导数都有一致极限；沿坐标线段积分可逐阶验证，导数极限正是函数极限的导数。各图中极限与已经得到的 \(u_*\) 相同，故相互相容，并说明 \(u_*\) 光滑。

	现在从 \(u_*\) 重新应用短时存在定理。旧解的全部空间导数收敛，而式\eqref{g3:eq:harmonic-heat-coordinate-time-recursion}及其逐次时间导数又给出全部时间导数在端点的极限；这些极限正是由初值 \(u_*\) 和同一方程确定的导数。故把新解接在 \(T_{\max}\) 后面，得到一个跨过该时刻的光滑解。这也是\cref{g3:thm:quasilinear-parabolic-continuation}在此处的具体实施：目标及一阶微分保持在紧集内，主部始终为固定的强抛物算子，所有高阶范数均受控制。延伸与最大性矛盾，故 \(T_{\max}=\infty\)。

	任意两个具有同一初值的光滑解在短时内相等；从每个已相等的时刻重新使用短时唯一性，便在任意有限时间区间上相等。各项先验界原已在任意光滑存在区间上成立，因此适用于所得全局解。
\end{proof}

这一证明把有限时间端点的问题化为一个具体的光滑极限问题。下一节则讨论另一种极限：沿 \(t\to\infty\) 取时刻，并借能量耗散判断所得映射为何调和。
% 纲要条目排为编号小节标题，并进入目录与书签；说明段仍保留于纲要注释中。

% !TeX root = ../../Book3.tex
% 纲要标题：### 26.5 调和代表的极限构造
% 纲要位置：几何学计划纲要.md，第 3230 行。
% 正文状态：初稿；尚未正式审核。

% 纲要内容开始
% 
% * 时间子列上的速度衰减
% * 统一高阶界及光滑子列
% * 极限满足调和方程
% * 同伦类的保持
% * 全时收敛与调和代表的唯一性
% 
% ---
% 
% 纲要内容结束

\section{调和代表的极限构造}
\label{g3:sec:2605}

热流在有限时间内保持光滑，只保证可以继续演化；要得到调和映射，还须在无穷时间中找到趋于静止的状态。能量恒等式提供时间方向的积分控制，高阶估计提供空间方向的紧性。两者结合，才能把一个演化过程变成同伦类中的调和代表。
\ProseParagraphBreak
以下源 \((M^m,g)\) 与目标 \((N^n,h)\) 均为闭连通光滑黎曼流形，\(m,n\ge1\)，且 \(\operatorname{Sect}_N\le0\)；一维目标的曲率条件理解为曲率张量为零。给定光滑 \(u_0:M\to N\)，以 \(u(x,t)\) 表示\cref{g3:thm:harmonic-heat-global-existence}给出的全时光滑热流，满足 \(\partial_tu=\tau(u)\)、\(u(\cdot,0)=u_0\)。空间上的协变导数同时使用固定的源联络与沿 \(u(\cdot,t)\) 的拉回联络，能量仍取 \(E(u)=\frac12\int_M|\mathrm du|^2\mu_g\)。

\subsection{时间子列上的速度衰减}
\label{g3:subsec:2605:01}
令
\[
	\mathcal E(t)=E(u(\cdot,t)),\qquad
	Q(t)=\int_M|\partial_tu(x,t)|^2\,\mu_g.
\]
\Cref{g3:thm:harmonic-heat-energy}给出 \(\mathcal E'(t)=-Q(t)\)。因此 \(\mathcal E\) 非增且非负，存在极限 \(\mathcal E_\infty\ge0\)，并有
\begin{equation}\label{g3:eq:harmonic-heat-energy-tail}
	\int_T^\infty Q(t)\,\mathrm dt
	=\mathcal E(T)-\mathcal E_\infty
	\longrightarrow0\qquad(T\to\infty).
\end{equation}
这个等式把剩余能量的下降量识别为后续时间中的总平方速度。

\begin{lemma}\label{g3:lem:harmonic-heat-slow-times}
	设 \((M^m,g)\)、\((N^n,h)\) 为闭连通光滑黎曼流形，\(m,n\ge1\)，\(u:M\times[0,\infty)\to N\) 为光滑调和映射热流，并记 \(Q(t)=\int_M|\partial_tu|^2\mu_g\)。存在 \(t_j\in[j,j+1]\)，使
	\[
		\|\partial_tu(\cdot,t_j)\|_{L^2(M)}\longrightarrow0.
	\]
	若 \(\operatorname{Sect}_N\le0\)，则进一步有 \(Q(t)\to0\) 及 \(tQ(t)\to0\)。
\end{lemma}
\begin{proof}
	在每个长度为一的区间上，连续非负函数 \(Q\) 至少有一点满足
	\[
		Q(t_j)\le\int_j^{j+1}Q(t)\,\mathrm dt.
	\]
	式\eqref{g3:eq:harmonic-heat-energy-tail}使右端趋于零，这就选出了所需时间。

	非正目标曲率下，\cref{g3:prop:harmonic-heat-speed}说明
	\((\partial_t+\Delta_g)|\partial_tu|^2\le0\)。在闭源上积分，\(\int_M\Delta_g f\,\mu_g=0\)，故 \(Q'(t)\le0\)。于是对 \(t>0\)，
	\[
		\frac t2Q(t)\le\int_{t/2}^tQ(s)\,\mathrm ds
		\le\mathcal E(t/2)-\mathcal E_\infty\longrightarrow0.
	\]
	这同时给出两项全时衰减。
\end{proof}

第一项只使用能量下降及下界，因而适用于闭源上任何已知全时存在的光滑热流。非正曲率还使平方速度的空间积分单调，遂将子列上的结论加强到全部充分大的时间。这里控制的是速度的平方积分；它与 \(\int_0^\infty\|\partial_tu\|_{L^2}\,\mathrm dt<\infty\) 是不同条件，后一个条件才直接给出 \(L^2\) 距离意义下的有限轨道长度。

\subsection{统一高阶界及光滑子列}
\label{g3:subsec:2605:02}
前节的\cref{g3:thm:harmonic-heat-uniform-derivatives}给出，对每个整数 \(r\ge0\)，存在与时间无关的 \(C_r\)，使
\begin{equation}\label{g3:eq:harmonic-heat-orbit-derivatives}
	\sup_{t\ge1}\sup_{x\in M}|\nabla^r\mathrm du(x,t)|\le C_r.
\end{equation}
这时源和目标都是固定流形；所需紧性只涉及同一对流形之间的映射。

\begin{lemma}\label{g3:lem:compact-target-map-smooth-precompactness}
	设 \((M,g)\)、\((N,h)\) 为闭连通光滑黎曼流形，维数均为正。若光滑映射列 \(v_j:M\to N\) 满足
	\[
		\sup_j\sup_M|\nabla^r\mathrm dv_j|<\infty
		\qquad\text{对每个整数 }r\ge0,
	\]
	则它有子列在 \(M\) 上按 \(C^\infty\) 收敛到一个光滑映射 \(v_\infty:M\to N\)。这里 \(C^\infty\) 收敛指在包含极限像的固定局部坐标中，各阶导数均一致收敛。
\end{lemma}
\begin{proof}
	先取得一致极限，再把每一阶导数的极限接上去。\(|\mathrm dv_j|\le C_0\) 与源中的最短测地线给出
	\[
		d_N(v_j(x),v_j(y))\le C_0d_M(x,y).
	\]
	目标紧，因而映射列等度连续且所有取值都留在紧集内。紧度量空间版本的 Arzel\`a--Ascoli 定理给出一致收敛子列，极限 \(v_\infty\) 连续并取值于 \(N\)。这一步也可直接从源的可数稠密点集上逐点选子列、再用上述统一连续模升级为一致收敛得到。

	对每个 \(x\in M\)，先选包含 \(v_\infty(x)\) 的目标坐标图 \(V\)，再选源坐标邻域 \(U\) 并缩小，使 \(v_\infty(\overline U)\) 紧含于 \(V\)。一致收敛保证充分大的 \(j\) 的 \(v_j(\overline U)\) 留在 \(V\) 的同一个紧子集中。源的紧性使这样的 \(U\) 只需有限个。

	在这些固定坐标中，坐标导数与协变导数由联络系数联系。例如
	\[
		\partial_i\partial_kv_j^\alpha
		=(\nabla\mathrm dv_j)_{ik}^{\alpha}
		+\Gamma^\ell_{ik}\partial_\ell v_j^\alpha
		-\Gamma^{\alpha}_{\beta\gamma}(v_j)
		\partial_i v_j^\beta\partial_k v_j^\gamma.
	\]
	源、目标坐标均限制在固定紧子集内，度量、逆度量及联络系数的各阶导数有界。上式先给二阶坐标导数界；重复求导，并利用给定的 \(\nabla^r\mathrm dv_j\) 界，归纳得到所有坐标导数的一致界。每个指定阶的导数又因下一阶有界而等度连续。

	依次对有限张坐标图中的零阶、一阶、二阶等坐标导数使用 Arzel\`a--Ascoli 定理，选取嵌套子列，再取对角子列。所得各阶一致极限确为 \(v_\infty\) 的相应导数：在较小坐标球中沿坐标线段使用
	\[
		v_j^\alpha(z+se_i)-v_j^\alpha(z)
		=\int_0^s\partial_i v_j^\alpha(z+re_i)\,\mathrm dr,
	\]
	便把一阶导数的极限与函数极限相联系；对高阶导数重复同一论证即可。坐标交叠处的极限由光滑换坐标与原映射的一致性相容，所以它们组成整体光滑映射，并得到 \(C^\infty\) 收敛。
\end{proof}

对任意 \(s_j\to\infty\)，将 \(v_j=u(\cdot,s_j)\) 代入该引理，就得到光滑收敛子列。特别地，上一小节选出的低速度时间列也有这样的子列；以后仍用 \(t_j\) 表示它。

\subsection{极限满足调和方程}
\label{g3:subsec:2605:03}
光滑收敛保留张力场，因为张力只使用映射的二阶导数以及固定源、目标度量的光滑系数。与速度衰减结合，就能判断所得聚点满足什么方程。

\begin{proposition}\label{g3:prop:harmonic-heat-limit-harmonic}
	设 \(M,N\) 为闭连通光滑黎曼流形，\(\operatorname{Sect}_N\le0\)，\(u\) 为由光滑初值产生的全时调和映射热流。若 \(t_j\to\infty\) 且 \(u(\cdot,t_j)\to v\) 于 \(C^\infty(M,N)\)，则 \(v\) 光滑调和，且
	\[
		E(v)=\lim_{t\to\infty}E(u(\cdot,t))=\mathcal E_\infty.
	\]
\end{proposition}
\begin{proof}
	在上一引理所选的固定坐标中，张力场为
	\[
		\tau(u)^\alpha
		=g^{ik}\left(\partial_i\partial_ku^\alpha
		-\Gamma^\ell_{ik}\partial_\ell u^\alpha
		+\Gamma^\alpha_{\beta\gamma}(u)
		\partial_i u^\beta\partial_k u^\gamma\right).
	\]
	映射的 \(C^2\) 收敛及目标系数的光滑性使这些分量一致趋于 \(\tau(v)^\alpha\)。目标内积的系数也一致收敛，因此
	\[
		\int_M|\tau(v)|^2\,\mu_g
		=\lim_{j\to\infty}\int_M|\tau(u(\cdot,t_j))|^2\,\mu_g
		=\lim_{j\to\infty}Q(t_j)=0.
	\]
	最后一步用\cref{g3:lem:harmonic-heat-slow-times}的全时速度衰减；若只取最初选定的时间子列，其子列结论已经足够。连续非负函数 \(|\tau(v)|^2\) 的积分为零，故它处处为零，得到 \(\tau(v)=0\)。同样，\(C^1\) 收敛使能量密度一致收敛；在紧源上积分便得能量等式。
\end{proof}

这一证明还说明，沿整个热流有
\begin{equation}\label{g3:eq:harmonic-heat-tension-decay}
	\sup_M|\nabla^r\tau(u(\cdot,t))|\longrightarrow0
	\qquad(t\to\infty),\quad r=0,1,2,\ldots.
\end{equation}
事实上，若对某个 \(r\) 不成立，就有 \(t_j\to\infty\) 和 \(\varepsilon>0\)，使左侧至少为 \(\varepsilon\)。由光滑预紧性选出 \(u(\cdot,t_j)\to v\) 的光滑子列；上一命题给出 \(\tau(v)=0\)，而张力的 \(r\) 阶协变导数是映射至 \(r+2\) 阶导数的光滑表达，其范数应趋于零，矛盾。

\subsection{同伦类的保持}
\label{g3:subsec:2605:04}
每个有限时间截面都与初值同伦，因为 \((x,s)\mapsto u(x,st)\)、\(0\le s\le1\)，就是一个同伦。还须把这一事实传到时间趋于无穷的极限。这里有用的是一致收敛与目标的正注入半径。

\begin{lemma}\label{g3:lem:close-maps-geodesic-homotopy}
	设 \((M,g),(N,h)\) 为闭光滑黎曼流形，\(N\) 连通，取 \(0<\rho<\operatorname{inj}_h(N)\)。若光滑映射 \(f_0,f_1:M\to N\) 满足
	\[
		\sup_{x\in M}d_N(f_0(x),f_1(x))<\rho,
	\]
	则二者光滑同伦，且一个同伦为
	\begin{equation}\label{g3:eq:close-map-geodesic-homotopy}
		H(x,s)=\exp_{f_0(x)}\bigl(s\log_{f_0(x)}f_1(x)\bigr),
		\qquad0\le s\le1.
	\end{equation}
\end{lemma}
\begin{proof}
	\Cref{g3:prop:compact-injectivity-lower-bound}保证可选共同的正 \(\rho\)，并说明距离小于 \(\rho\) 的两点由唯一长度小于 \(\rho\) 的最短测地线连接。为核对公式的光滑性，考虑
	\[
		\mathcal F:\{(p,w)\in\mathrm TN:|w|<\rho\}\longrightarrow N\times N,
		\qquad\mathcal F(p,w)=(p,\exp_pw).
	\]
	其像正是 \(d_N(p,q)<\rho\) 的点对，且它在各纤维上单射。微分的第一分量先确定底点方向，第二分量中关于 \(w\) 的微分可逆，故整体微分可逆。由逆函数定理，\(\mathcal F\) 的逆在像上光滑，其第二分量就是 \(\log_pq\)。式\eqref{g3:eq:close-map-geodesic-homotopy}于是光滑，且在 \(s=0,1\) 分别等于 \(f_0,f_1\)。
\end{proof}

同伦调和代表的极限构造与同伦类内能量极小者的存在性，见\cite[第11.A节定理及推论，第156--158页]{EellsSampson1964HarmonicMaps}。本章以能量空间中的短时理论和内在先验估计证明紧目标版本，并通过光滑逼近涵盖连续映射。

\termidx[Eells Sampson dingli]{Eells--Sampson 定理}
\foreignidx{Eells--Sampson theorem}
\begin{theorem}[Eells--Sampson 型调和代表]\label{g3:thm:harmonic-heat-homotopy-representative}
	设 \((M^m,g)\)、\((N^n,h)\) 为闭连通光滑黎曼流形，\(m,n\ge1\)，且 \(\operatorname{Sect}_N\le0\)。则每个光滑映射 \(u_0:M\to N\) 的调和映射热流都全时光滑存在，并有时间子列 \(t_j\to\infty\)，使 \(u(\cdot,t_j)\) 在 \(C^\infty\) 中收敛到与 \(u_0\) 同伦的光滑调和映射 \(u_\infty\)。特别地，每个连续映射 \(M\to N\) 都同伦于一个光滑调和映射。
\end{theorem}
\begin{proof}
	全时存在由\cref{g3:thm:harmonic-heat-global-existence}给出。\Cref{g3:lem:harmonic-heat-slow-times}选出 \(L^2\) 速度趋于零的时间，式\eqref{g3:eq:harmonic-heat-orbit-derivatives}与\cref{g3:lem:compact-target-map-smooth-precompactness}再选光滑收敛子列。\Cref{g3:prop:harmonic-heat-limit-harmonic}说明极限 \(u_\infty\) 调和。

	取 \(0<\rho<\operatorname{inj}_h(N)\)。一致收敛使充分大的 \(j\) 满足
	\(\sup_Md_N(u(x,t_j),u_\infty(x))<\rho\)。前一引理将 \(u(\cdot,t_j)\) 与 \(u_\infty\) 用短测地线同伦连接，再与 \(u_0\) 到 \(u(\cdot,t_j)\) 的热流同伦拼接，就得 \(u_0\simeq u_\infty\)。拼接后的连续同伦已经足够；若要求光滑，可在两个时间段的端点用光滑且在端点附近恒定的参数函数重参数化。

	最后，第二卷第14.1.2小节关于连续映射与光滑同伦的推论已经证明：光滑流形间每个连续映射都同伦于光滑映射，源允许边界，目标无边界。将它用于这里的闭 \(M,N\)，先把给定连续映射换成光滑的 \(u_0\)，再用刚才的构造即可。
	% 跨卷依赖：g2:cor:continuous-smooth-homotopy。
\end{proof}

第十八章的球面气泡例子中，弱收敛可以丢失度数。这里的一致收敛使每个点的末段运动都局限在同一个可光滑选取的短测地线内，因而保留了同伦类；高阶界同时保证张力场与能量都能传到极限。

\subsection{全时收敛与调和代表的唯一性}
\label{g3:subsec:2605:05}
上一定理先得到聚点。若同一轨道的所有可能聚点都被唯一性确定，光滑预紧性就会把子列收敛升级为全时收敛。

\begin{proposition}\label{g3:prop:harmonic-heat-unique-limit}
	设 \(M,N\) 为闭连通光滑黎曼流形，\(\operatorname{Sect}_N\le0\)。若光滑映射 \(u_0:M\to N\) 的同伦类中恰有一个光滑调和映射 \(v\)，则从 \(u_0\) 出发的调和映射热流满足
	\[
		u(\cdot,t)\longrightarrow v\quad\text{于 }C^\infty(M,N),
		\qquad t\to\infty.
	\]
\end{proposition}
\begin{proof}
	任取 \(t_j\to\infty\)，光滑预紧性给出一个光滑收敛子列。其极限由\cref{g3:prop:harmonic-heat-limit-harmonic}调和，再由短测地线同伦与 \(u_0\) 同伦，故只能是 \(v\)。

	若全时收敛不成立，就有 \(C^\infty(M,N)\) 中 \(v\) 的一个开邻域 \(\mathcal U\)，以及 \(s_j\to\infty\)，使 \(u(\cdot,s_j)\notin\mathcal U\)。对该时间列再取光滑收敛子列，其极限仍为 \(v\)，因而最终应落入 \(\mathcal U\)，矛盾。这同时证明每个有限阶的收敛。
\end{proof}

为了给出可核验的唯一性条件，需要比较同伦类中任意两个调和映射。目标本身可能有多条连接同端点的测地线；在它的万有覆盖上，Cartan--Hadamard 定理给出唯一测地线，并使这个比较成为全局构造。

\begin{lemma}\label{g3:lem:nonpositive-target-geodesic-homotopy}
	设 \((M,g)\) 为闭连通光滑黎曼流形，\((N,h)\) 为连通完备光滑无边界黎曼流形，维数均为正，且 \(\operatorname{Sect}_N\le0\)。任意两个同伦的光滑映射 \(f_0,f_1:M\to N\) 之间存在光滑同伦 \(F(x,s)\)，使每条 \(s\mapsto F(x,s)\) 都是以 \([0,1]\) 为仿射参数区间的测地线。令 \(F_s(x)=F(x,s)\)、\(V=\partial_sF\)，则
	\begin{equation}\label{g3:eq:harmonic-geodesic-homotopy-convexity}
		\frac{\mathrm d^2}{\mathrm ds^2}E(F_s)
		=\int_M\left(|D V|^2
		-\sum_i\langle R^N(V,\mathrm dF_s(e_i))\mathrm dF_s(e_i),V\rangle\right)\mu_g
		\ge0.
	\end{equation}
\end{lemma}
\begin{proof}
	取光滑万有覆盖 \(p:\widetilde M\to M\)、\(q:\widetilde N\to N\)，目标覆盖带拉回度量。\Cref{g3:prop:riemannian-universal-cover}保证 \(\widetilde N\) 完备、单连通，覆盖变换等距；曲率仍非正，故它是 Hadamard 流形。

	设 \(H:M\times[0,1]\to N\) 为给定的连续同伦。因为 \(\widetilde M\times[0,1]\) 单连通，覆盖提升判据给出 \(H\circ(p\times\operatorname{id})\) 的一个提升 \(\widetilde H\)。这里使用第一卷第5.3.3小节的提升判据：单连通源的基本群像为平凡群，因而包含于覆盖所对应的子群。记两端为 \(\widetilde f_0,\widetilde f_1\)；原映射光滑且覆盖局部微分同胚，所以这两个提升光滑。
	% 跨卷依赖：g1:thm:lifting-criterion。

	对源的每个覆盖变换 \(\gamma\)，存在目标覆盖变换 \(\rho(\gamma)\)，使
	\[
		\widetilde H(\gamma\widetilde x,s)
		=\rho(\gamma)\widetilde H(\widetilde x,s).
	\]
	确实，在一个指定点先由目标覆盖变换在纤维上的传递性选择 \(\rho(\gamma)\)，两边随后是同一个映射的提升，在连通域中由提升唯一性处处相等。自由性还保证该变换唯一，并由复合可知 \(\rho\) 是群同态。因此两个端点提升具有相同的等变关系。

	由\cref{g3:prop:hadamard-geodesic-uniqueness}，可定义
	\[
		\widetilde F(\widetilde x,s)
		=\exp_{\widetilde f_0(\widetilde x)}
		\left(s\log_{\widetilde f_0(\widetilde x)}
		\widetilde f_1(\widetilde x)\right).
	\]
	测地线随两端点光滑变化，所以 \(\widetilde F\) 光滑。目标覆盖变换等距，保持测地线及仿射参数；唯一性使 \(\widetilde F\) 具有同一等变关系。故 \(q\circ\widetilde F\) 对源的覆盖变换不变，下降为所求 \(F\)。覆盖局部等距保证下降后的各条参数曲线仍是测地线，因而 \(D_sV=0\)。

	最后用\cref{g3:thm:harmonic-map-first-variation,g3:prop:harmonic-map-jacobi-linearization}对能量连续求两次导数：
	\[
		\frac{\mathrm d^2}{\mathrm ds^2}E(F_s)
		=-\int_M\bigl(\langle D_s\tau(F_s),V\rangle
		+\langle\tau(F_s),D_sV\rangle\bigr)\mu_g
		=\int_M\langle J_{F_s}V,V\rangle\,\mu_g.
	\]
	这里 \(D_sV=0\)，故无需假设中间的映射 \(F_s\) 调和。闭源上的分部积分式\eqref{g3:eq:harmonic-map-jacobi-green}给出式\eqref{g3:eq:harmonic-geodesic-homotopy-convexity}。截面曲率非正使每个被减去的曲率项非正，因而右端非负。
\end{proof}

\begin{corollary}\label{g3:cor:nonpositive-harmonic-homotopy-minimum}
	设 \(M\) 为闭连通光滑黎曼流形，\(N\) 为连通完备光滑无边界黎曼流形，维数均为正，且 \(\operatorname{Sect}_N\le0\)。以 \(E(v)=\tfrac12\int_M|\mathrm dv|^2\,\mathrm d\mu_g\) 记能量。则光滑调和映射 \(f:M\to N\) 在所有与它同伦的光滑映射中使能量最小。
\end{corollary}
\begin{proof}
	对任意同伦的光滑映射 \(v\)，取前一引理的测地同伦，并令 \(A(s)=E(F_s)\)。调和性给出 \(A'(0)=0\)，凸性给出 \(A'\) 非减，故
	\(E(v)-E(f)=\int_0^1A'(s)\,\mathrm ds\ge0\)。
\end{proof}

\begin{theorem}\label{g3:thm:negative-target-harmonic-uniqueness}
	设 \((M^m,g)\) 为闭连通光滑黎曼流形，\((N^n,h)\) 为连通完备光滑无边界黎曼流形，\(m,n\ge2\)，且 \(\operatorname{Sect}_N<0\)。若光滑调和映射 \(f_0:M\to N\) 在某一点 \(x_0\) 满足 \(\operatorname{rank}(\mathrm df_0)_{x_0}\ge2\)，则其同伦类中只有这一个光滑调和映射。若 \(N\) 还闭，则从该同伦类中任意光滑初值出发的调和映射热流都在 \(C^\infty\) 中全时收敛到 \(f_0\)。
\end{theorem}
\begin{proof}
	设 \(f_1\) 是同伦的另一个调和映射，取上述测地同伦。令 \(A(s)=E(F_s)\)，则 \(A''\ge0\)，而两端的调和性使 \(A'(0)=A'(1)=0\)。非减函数 \(A'\) 因此恒为零，故式\eqref{g3:eq:harmonic-geodesic-homotopy-convexity}的右端对每个 \(s\) 都等于零。

	特别在 \(s=0\)，各项均非负且连续，因而
	\[
		DV_0=0,\qquad
		\langle R^N(V_0,\mathrm df_0(e_i))\mathrm df_0(e_i),V_0\rangle=0
	\]
	处处成立。严格负曲率下，第二个等式说明 \(V_0(x)\) 与每个 \(\mathrm df_0(e_i)\) 线性相关，除非 \(V_0(x)=0\)。在 \(x_0\) 有至少两个线性无关的像方向，所以只能有 \(V_0(x_0)=0\)。又因 \(DV_0=0\)，其长度在连通源上恒定，故 \(V_0\equiv0\)。每条 \(s\mapsto F(x,s)\) 的初速度于是为零；测地初值唯一性使该曲线恒定，得到 \(f_1=f_0\)。若 \(N\) 闭，再用\cref{g3:prop:harmonic-heat-unique-limit}即得全时收敛。
\end{proof}

秩条件排除了常值映射和沿同一闭测地线移动的退化情形。非正曲率允许同伦类中存在整个调和映射族；热流仍可能按初值选定其中的一个。下面可以直接看清这个选择。

\begin{example}\label{g3:exm:flat-circle-harmonic-heat-flow}
	把源与目标都取为 \(S^1=\mathbb R/(2\pi\mathbb Z)\)，度量为 \(\mathrm d\theta^2\)。一个光滑映射的提升可写成
	\[
		\widetilde u_0(\theta)=k\theta+b(\theta),
		\qquad k\in\mathbb Z,\quad b(\theta+2\pi)=b(\theta).
	\]
	这里提升满足 \(\widetilde u_0(\theta+2\pi)-\widetilde u_0(\theta)=2\pi k\)：差除以 \(2\pi\) 是连续整数值函数，故恒定。整数 \(k\) 就是映射的度数。若
	\[
		b(\theta)=\sum_{\ell\in\mathbb Z}\widehat b_\ell e^{i\ell\theta},
		\qquad
		\widehat b_\ell=\frac1{2\pi}\int_0^{2\pi}b(\theta)e^{-i\ell\theta}\,\mathrm d\theta,
	\]
	则热流由下式给出：
	\begin{equation}\label{g3:eq:circle-harmonic-heat-fourier}
		u([\theta],t)=\left[k\theta+\widehat b_0
		+\sum_{\ell\ne0}e^{-\ell^2t}\widehat b_\ell e^{i\ell\theta}\right].
	\end{equation}
	因为 \(b\) 光滑，其 Fourier 系数快速衰减，级数及任意指定阶的空间、时间导数均可逐项求导，并在 \(t\ge0\) 一致收敛。目标在角坐标中平坦，张力为 \(\partial_\theta^2\widetilde u\)，故上式满足 \(\partial_tu=\tau(u)\)，在 \(t=0\) 等于初值。提升增加 \(2\pi k\) 的等变关系始终保持，所以它确实定义圆周之间的热流。
	\[
		u([\theta],t)\longrightarrow[k\theta+\widehat b_0]
		\quad\text{于 }C^\infty,\qquad t\to\infty.
	\]
	对每个 \(r\)，非零频率部分的 \(C^r\) 范数至多为
	\(e^{-t}\sum_{\ell\ne0}|\ell|^r|\widehat b_\ell|\)，所以收敛还是指数的。Parseval 等式则给出
	\[
		E(u(t))=\pi k^2
		+\pi\sum_{\ell\ne0}\ell^2|\widehat b_\ell|^2e^{-2\ell^2t},
		\qquad
		Q(t)=2\pi\sum_{\ell\ne0}\ell^4|\widehat b_\ell|^2e^{-2\ell^2t},
	\]
	可逐项核对 \(E'=-Q\)。
	同一度数类中的所有 \([\theta]\mapsto[k\theta+c]\) 都调和，且随 \(c\) 连续变化，彼此同伦；当 \(k=0\) 时就是所有常值映射。极限中的 \([\widehat b_0]\) 由初值决定：改换提升只把它增加 \(2\pi\) 的整数倍，不改变圆周值。因而这一模型既有全时收敛，也有同伦类中调和代表的非唯一性。
\end{example}
% 纲要条目排为编号小节标题，并进入目录与书签；说明段仍保留于纲要注释中。


\begin{chaptersummary}
	抛物能量空间 \(X_T^s=C_tH^s\cap L^2_tH^{s+1}\)、\(u_t\in L^2_tH^{s-1}\) 同时记录连续初值、空间平滑和时间变化。强抛物主符号经冻结系数给出 Gårding 估计；与固定正椭圆算子交换的谱投影则使有限维逼近保留这一能量控制。所得线性初值同构的逆范数在短区间上一致有界，正是拟线性构造所需的输入。

	\ProseParagraphBreak
	拟线性方程在时间常值初值附近的误差具有 \(\sqrt T\) 因子。先选非线性邻域，再缩短时间，压缩映射给出解；参数隐函数定理与源上的光滑坐标变化共同恢复全部正则性。光滑解的差满足强抛物线性方程，故解唯一。若一阶喷射留在强抛物紧区域且二阶导数有界，乘积估计和高阶能量不等式保证全部高阶范数有界，进而可在有限端点重新求解。

	\ProseParagraphBreak
	调和映射热流满足 \(E'(t)=-\|u_t\|_2^2\)。非正目标曲率使能量密度成为带源 Ricci 零阶项的次解，并使平方速度成为次热解。能量的 \(L^1\) 界通过每次长度相同的时间窗转成统一梯度界；导数演化中的负最高阶项再经加权最大值估计给出全阶控制。这些界排除有限时奇点，也使大时间轨道光滑预紧。

	\ProseParagraphBreak
	能量耗散选出低速度时间，光滑紧性给出调和极限，目标的正注入半径则用短测地线同伦保留初始同伦类。非正曲率下的测地同伦使能量凸，因此调和代表在同伦类中能量最小。若代表唯一，预紧性将子列收敛升级为全时收敛；严格负曲率加上一点处秩至少为二，是本章证明的一个具体唯一性条件。平坦圆周的显式流说明，同伦类中代表不唯一时，特定初值的热流仍可收敛。
\end{chaptersummary}

\BookExercises

\begin{exercise}\label{g3:ex:heat-parabolic-scaling}
	设 \(u\) 是 \(\mathbb R^n\times[0,T)\) 上的光滑热方程解，\(u_t=\sum_i u_{ii}\)，\(\rho>0\)。证明 \(u_\rho(x,t)=u(\rho x,\rho^2t)\) 仍是解。若初值为 \(u_0(x)=\cos(kx_1)\)，求一个有界解并说明频率 \(k\) 的自然衰减时间为何与 \(k^{-2}\) 同阶。
\end{exercise}
\begin{solution}
	链式法则给 \(\partial_tu_\rho=\rho^2u_t\)、\(\sum_i\partial_i^2u_\rho=\rho^2\sum_i u_{ii}\)，所以方程不变，存在时间变成 \(T/\rho^2\)。一个有界解为 \(e^{-k^2t}\cos(kx_1)\)。当 \(k\ne0\) 时，振幅衰减到初值的 \(e^{-1}\) 倍恰需时间 \(k^{-2}\)；空间长度缩小为原来的一半，对应的时间尺度缩小为原来的四分之一。\(k=0\) 是不衰减的常数模式。
\end{solution}

\begin{exercise}\label{g3:ex:backward-heat-discontinuity}
	在标准圆周上考虑 \(v_t-\Delta v=0\)，其中 \(\Delta=-\partial_x^2\)。取初值 \(v_m(0,x)=e^{-m}\cos(mx)\)，\(m=1,2,\ldots\)。证明初值在 \(C^\infty\) 中趋于零，但任意固定正时间上的解不趋于零。说明这与强抛物条件的关系。
\end{exercise}
\begin{solution}
	任意固定 \(r\) 阶导数的初始上界为 \(m^re^{-m}\to0\)。方程的解却是 \(v_m(t,x)=e^{m^2t-m}\cos(mx)\)，其上确界在每个 \(t>0\) 趋于无穷。因而求解映射即使在这些光滑初值上也不连续。把方程写成 \(v_t=-v_{xx}\)，二阶系数是 \(-1\)，不满足正的强抛物条件；时间反向使原本衰减的高频成分迅速放大。
\end{solution}

\begin{exercise}\label{g3:ex:heat-periodic-forcing}
	设闭流形上 \(\Delta\phi=\lambda\phi\)，\(\lambda>0\)，且 \(\omega>0\)。求零初值问题 \(u_t+\Delta u=\cos(\omega t)\phi\) 的解，说明它的长时间行为。另求空间常值外力 \(\cos(\omega t)\) 对应的零初值解。
\end{exercise}
\begin{solution}
	令 \(u(t)=a(t)\phi\)，则 \(a'+\lambda a=\cos(\omega t)\)、\(a(0)=0\)。直接代入得
	\[
		a(t)=\frac{\lambda\cos(\omega t)+\omega\sin(\omega t)-\lambda e^{-\lambda t}}
		{\lambda^2+\omega^2}.
	\]
	解趋近周期响应 \((\lambda\cos(\omega t)+\omega\sin(\omega t))\phi/(\lambda^2+\omega^2)\)，误差指数衰减，但解本身不趋向驻定状态。常值模式满足 \(a'=\cos(\omega t)\)，所以 \(u(t,x)=\sin(\omega t)/\omega\)。此例说明持续外力的时间行为会传入长期解，热平滑本身不能保证驻定极限。
\end{solution}

\begin{exercise}\label{g3:ex:heat-energy-space-sharpness}
	在标准圆周上，固定整数 \(s\ge0\) 和 \(T>0\)，令 \(u_m(t,x)=m^{-s}e^{-m^2t}\cos(mx)\)。证明 \(\|u_m\|_{X_T^s}\) 对 \(m\ge1\) 一致有界，而 \(\sup_{0\le t\le T}\|u_m(t)\|_{H^{s+1}}\) 无界。所有 Sobolev 范数可用 \(I+\Delta\) 的谱范数。
\end{exercise}
\begin{solution}
	圆周上 \(\|\cos(mx)\|_2^2=\pi\)，故
	\[
		\sup_t\|u_m(t)\|_{H^s}^2=\pi(1+m^2)^sm^{-2s}\le\pi2^s.
	\]
	又 \(\int_0^Te^{-2m^2t}\,\mathrm dt\le(2m^2)^{-1}\)，从而
	\[
		\int_0^T\|u_m\|_{H^{s+1}}^2\,\mathrm dt
		\le\frac\pi2(1+m^2)^{s+1}m^{-2s-2},
	\]
	以及
	\[
		\int_0^T\|(u_m)_t\|_{H^{s-1}}^2\,\mathrm dt
		\le\frac\pi2(1+m^2)^{s-1}m^{2-2s}.
	\]
	两式均一致有界，包括 \(s=0\) 的负一阶范数。然而初时刻的 \(H^{s+1}\) 范数平方为 \(\pi(1+m^2)^{s+1}m^{-2s}\)，与 \(m^2\) 同阶。额外一阶空间导数只能按时间平方积分控制，不能把线性估计中的 \(L_t^2H^{s+1}\) 任意换成 \(C_tH^{s+1}\)。
\end{solution}

\begin{exercise}\label{g3:ex:reaction-diffusion-lifespan}
	设 \((M,g)\) 闭，光滑初值满足 \(u_0\ge a>0\)。证明 \(u_t+\Delta_gu=u^2\) 的光滑解不能存在超过时间 \(a^{-1}\)。
\end{exercise}
\begin{solution}
	令 \(b(t)=a/(1-at)\)，定义于 \(t<a^{-1}\)。在解的任意共同紧时间区间上，\(w=u-b\) 满足
	\[
		w_t+\Delta_gw-(u+b)w=0,\qquad w(0)\ge0.
	\]
	系数 \(u+b\) 在该区间有界，标量最大值原理作用于 \(-w\) 给出 \(u\ge b\)。如果解存在于一个包含 \(a^{-1}\) 的更长光滑区间，它在 \([0,a^{-1}]\times M\) 上应有界；但 \(b(t)\to\infty\)，矛盾。此处失控的是解的零阶大小，主部仍始终强抛物。
\end{solution}

\begin{exercise}\label{g3:ex:harmonic-heat-metric-scalings}
	设 \((M^m,g)\) 为闭光滑黎曼流形，\((N,h)\) 为光滑无边界黎曼流形，\(u:[0,T)\times M\to N\) 为光滑调和映射热流，将源、目标度量分别换为 \(cg\)、\(dh\)，其中 \(c,d>0\) 为常数。求能量与张力的变化，写出由原热流 \(u(t)\) 得到的新热流，并核对能量耗散公式的尺度。
\end{exercise}
\begin{solution}
	常数缩放不改变两个 Levi--Civita 联络。张力只对源指标取迹，因此 \(\tau_{cg,dh}=c^{-1}\tau_{g,h}\)。目标度量的缩放改变长度，但不改变这个取值于目标切空间的向量。逆源度量贡献 \(c^{-1}\)，源体积贡献 \(c^{m/2}\)，目标内积贡献 \(d\)，故
	\[
		E_{cg,dh}(f)=d\,c^{m/2-1}E_{g,h}(f).
	\]
	新流为 \(\widetilde u(t)=u(t/c)\)。其能量时间导数为 \(dc^{m/2-2}E'_{g,h}(u(t/c))\)，而新速度平方积分也等于 \(dc^{m/2-2}\|u_t(t/c)\|_{L^2(g,h)}^2\)。两边因子相同，耗散恒等式保持。
\end{solution}

\begin{exercise}\label{g3:ex:harmonic-heat-geodesic-image}
	设 \((M,g)\) 闭且连通，\(\gamma:\mathbb R\to N\) 是完备目标上的单位速测地线，\(v_0:M\to\mathbb R\) 光滑。以 \(T_t=e^{-t\Delta_g}\)、\(\Delta_g=-\operatorname{div}_g\nabla\) 记源热半群。证明 \(u(t,x)=\gamma(T_tv_0(x))\) 是调和映射热流，并证明它光滑收敛到常值映射 \(\gamma(\overline v_0)\)，其中 \(\overline v_0\) 是源上的平均值。
\end{exercise}
\begin{solution}
	令 \(v=T_tv_0\)。复合映射的张力公式在这里直接给出
	\[
		\tau(\gamma\circ v)=-\gamma'(v)\Delta_gv
		+(\nabla_{\gamma'}\gamma')(v)|\mathrm dv|^2
		=-\gamma'(v)\Delta_gv.
	\]
	而 \(u_t=\gamma'(v)v_t=-\gamma'(v)\Delta_gv\)，所以满足热流方程。连通闭源的热半群将 \(v_0-\overline v_0\) 在每个固定阶范数中指数衰减；最大值原理还使 \(v\) 留在 \([\min v_0,\max v_0]\)。\(\gamma\) 在这个紧区间上的各阶导数有界，链式法则遂给 \(u\to\gamma(\overline v_0)\) 的光滑收敛。该构造不需要目标曲率非正，说明非正曲率是控制任意初值的充分几何条件，某些特殊解可以具有更强的约束。
\end{solution}

\begin{exercise}\label{g3:ex:positive-source-heat-constant-limit}
	设闭连通源 \((M^m,g)\)、\(m\ge2\) 满足 \(\operatorname{Ric}_g\ge\kappa g\)、\(\kappa>0\)，紧连通目标 \((N,h)\) 满足 \(\operatorname{Sect}_N\le0\)。证明任意光滑初值的调和映射热流全时光滑收敛到一个常值映射。要求从演化方程给出一致距离意义下的指数收敛，并说明所得同伦结论。
\end{exercise}
\begin{solution}
	Bochner 公式给 \((\partial_t+\Delta_g)e\le-2\kappa e\)，故 \(\sup e(t)\le e^{-2\kappa t}\sup e(0)\)，并有 \(E(t)\le e^{-2\kappa t}E(0)\)。令 \(Q(t)=\int|u_t|^2\)。速度不等式使 \(Q\) 非增，所以对 \(t\ge1\)，
	\[
		Q(t)\le\int_{t-1}^tQ(s)\,\mathrm ds
		\le E(t-1)\le E(0)e^{-2\kappa(t-1)}.
	\]
	再把平方速度的次热方程在长度为一的时间窗内与热核比较。若 \(K_1=\max_{x,y}K(1,x,y)\)，则对 \(t\ge2\)，
	\[
		\sup_M|u_t(t)|^2\le K_1 Q(t-1)\le C e^{-2\kappa t}.
	\]
	从 \(t\) 积分到 \(b>t\)，得 \(\sup_xd_N(u(x,t),u(x,b))\le C'e^{-\kappa t}\)。目标完备使热流有一致极限 \(u_\infty\)，同一不等式给指数距离估计。又 \(\sup|\mathrm du(t)|\to0\)，源中最短测地线的长度估计说明极限在连通源上恒定。光滑预紧性与一致极限的唯一性使整个轨道光滑收敛。有限时间热流同伦加上末段短测地线同伦说明 \(u_0\) 零同伦；连续初值再先用光滑逼近。因此在这些几何条件下，每个连续映射 \(M\to N\) 都零同伦。
\end{solution}

\begin{exercise}\label{g3:ex:harmonic-representatives-convex-family}
	设 \(M\) 闭且连通，\(N\) 完备、连通且截面曲率非正。证明两个同伦的光滑调和映射具有相同能量，而且连接它们的测地同伦中的每个映射都调和。进一步证明零同伦的光滑调和映射只能是常值映射。
\end{exercise}
\begin{solution}
	取\cref{g3:lem:nonpositive-target-geodesic-homotopy}的同伦 \(F_s\)，令 \(A(s)=E(F_s)\)。能量凸性及两端调和给出 \(A''\ge0\)、\(A'(0)=A'(1)=0\)，所以 \(A'\equiv0\)，各能量相同。\Cref{g3:cor:nonpositive-harmonic-homotopy-minimum}说明初端已经在同伦类中能量最小，因而每个 \(F_s\) 也取得这个最小值。对 \(F_s\) 的任意光滑小变分，变分后的映射仍在该同伦类内，所以能量一阶导数为零。第一变分公式及任意变分场给出 \(\tau(F_s)=0\)。如果原映射零同伦，就与某个常值映射比较，得到其能量等于零。于是 \(\mathrm df=0\)，连通性给出 \(f\) 常值。
\end{solution}

\begin{exercise}\label{g3:ex:square-integrable-speed-not-convergence}
	在标准圆周上考虑曲线 \(\gamma(t)=[\log(1+t)]\)、\(t\ge0\)。证明它的平方速度在全时间可积，但曲线不收敛；构造一个非负函数 \(E(t)\)，使 \(E'(t)=-|\gamma'(t)|^2\)。说明这个例子能对热流极限证明提供什么提醒。
\end{exercise}
\begin{solution}
	\(|\gamma'(t)|=(1+t)^{-1}\)，故 \(\int_0^\infty|\gamma'|^2\,\mathrm dt=1\)，但 \(\int_0^\infty|\gamma'|\,\mathrm dt=\infty\)。取 \(t_j=e^{2\pi j}-1\)、\(s_j=e^{2\pi j+\pi}-1\)，两列像分别为 \([0]\)、\([\pi]\)，所以不收敛。函数 \(E(t)=(1+t)^{-1}\) 满足所需耗散式。该曲线不是这里的调和映射热流；它说明仅有非负能量与平方速度耗散，不能推出轨道收敛。热流论证还需要空间紧性以及识别全部聚点的条件。
\end{solution}

\begin{exercise}\label{g3:ex:negative-target-rank-one-family}
	设负截面曲率目标 \(N\) 有一条单位速闭测地线 \(\gamma:\mathbb R/L\mathbb Z\to N\)，\(L>0\)。对非零整数 \(k\) 与 \(a\in\mathbb R/L\mathbb Z\)，定义
	\[
		f_a:\mathbb R/(2\pi\mathbb Z)\to N,\qquad
		f_a([\theta])=\gamma\left(\left[\frac{kL}{2\pi}\theta+a\right]\right).
	\]
	证明这些映射调和、彼此同伦，计算能量，并说明严格负曲率唯一性定理中的秩条件为何必要。
\end{exercise}
\begin{solution}
	\(\theta\) 增加 \(2\pi\) 时参数增加 \(kL\)，故映射良定义。其速度为 \(\frac{kL}{2\pi}\gamma'\)，协变加速度为零，所以张力为零。沿参数 \(a\) 的任意连接道路作平移便得到同伦。速度长度恒为 \(|k|L/(2\pi)\)，因此
	\[
		E(f_a)=\frac12\,(2\pi)\left(\frac{kL}{2\pi}\right)^2
		=\frac{k^2L^2}{4\pi}.
	\]
	小的非零参数移动给出不同映射，因为 \(\gamma'\ne0\)。它们的秩处处为一，所以负曲率并没有排除沿同一测地线平移的自由；一点处秩至少为二正用于消除这一情形。
\end{solution}

\begin{exercise}\label{g3:ex:noncompact-target-escape}
	设 \(M\) 为非空闭流形，取 \(v_j:M\to\mathbb R\)、\(v_j\equiv j\)。核对它们的能量和全部正阶导数，说明为什么只有导数界仍不足以在非紧目标中抽取光滑收敛子列。给出一个可以代替目标紧性的附加条件。
\end{exercise}
\begin{solution}
	每个 \(v_j\) 都是调和映射，能量为零，且全部正阶导数恒为零。但是在任一点 \(x\in M\)，数列 \(v_j(x)=j\) 没有收敛子列，故映射列连逐点收敛子列都没有。失控的是取值位置。若所有映射的像共同包含在目标的一个固定紧集内，并有各阶协变导数的一致界，则\cref{g3:lem:compact-target-map-smooth-precompactness}的有限坐标及对角取列证明仍适用。也可在完备目标中要求一个基点的像有界，再结合统一的一阶界和源的直径；Hopf--Rinow 定理使包含全部像的闭有界球成为紧集。后一说法对连通源成立，若源不连通，则在各个分支分别选基点。
\end{solution}
